\documentclass[12pt,a4paper]{report} % or your university's class

% ---------- Encoding & language ----------
\usepackage[utf8]{inputenc}
\usepackage[T1]{fontenc}
\usepackage[english]{babel}

% ---------- Typography & layout ----------
\usepackage{lmodern}
\usepackage{setspace}
\onehalfspacing
\usepackage{geometry}
\geometry{margin=2.5cm}

% ---------- Math / equations ----------
\usepackage{amsmath,amssymb,amsthm,mathtools}

% ---------- Tables ----------
\usepackage{booktabs}
\usepackage{multirow}
\usepackage{multicol}
\usepackage{tabularx}
\usepackage{longtable}
\usepackage{siunitx}

% ---------- Figures ----------
\usepackage{graphicx}
\usepackage{float}
\usepackage{subcaption}
\usepackage{caption}

% ---------- References / hyperlinks ----------
\usepackage[hidelinks]{hyperref}
\usepackage{cleveref}
\usepackage{csquotes}

% ---------- Bibliography (biblatex + biber) ----------
\usepackage[
    backend=biber,
    style=ieee,      % Uses numeric citations [1] and IEEE bibliography format
    sorting=none     % CRITICAL: Sorts by order of appearance (Standard IEEE)
]{biblatex}
\addbibresource{reference.bib}

% ---------- Other useful packages ----------
\usepackage{enumitem}
\usepackage{xcolor}

\begin{document}

% ---------- Front matter ----------
\title{Trustworthy RAG: An Evaluation Agent for Detecting Misinformation and Knowledge Poisoning in Generative AI Systems}
\author{Balkrishna Giri}
\date{\today} % or fixed date

\maketitle

\begin{abstract}
Retrieval-Augmented Generation (RAG) systems improve large language model reliability by grounding responses in external knowledge. However, RAG architectures inherently trust their retrieval output, creating a critical Security-Reliability Gap: high semantic relevance does not guarantee factual truth. Adversaries can exploit this by injecting malicious documents into the knowledge corpus --- a threat known as Knowledge Poisoning --- causing the system to generate targeted misinformation.

This thesis proposes an autonomous Evaluation Agent as a defensive middleware within the RAG pipeline. The agent combines Natural Language Inference (NLI) for factual verification, a multi-signal Poison Detector (linguistic, structural, intra-document, cross-document, and semantic-outlier signals), and a composite Trust Index ($T = 0.4 \cdot F + 0.35 \cdot C + 0.25 \cdot (1-P)$) to produce an interpretable trustworthiness score for each generated response.

Experiments on TruthfulQA and FEVER benchmarks with four adversarial poisoning strategies (contradiction, instruction injection, entity swap, and subtle manipulation) demonstrate that the system achieves 91\% accuracy and 57.1\% F1 score with 100\% precision (zero false positives) using Llama~3.3~70B, a +7\% improvement over the naive always-trust baseline. Instruction injection is detected with 100\% recall, while entity-swap and subtle manipulation remain near-undetectable at the textual level, representing an architectural limit of surface-signal-based approaches.

A 2×2 factorial experiment (\{Llama~3.3~70B, Qwen~3.5~35B\} × \{all-MiniLM-L6-v2, Snowflake Arctic Embed2\}) reveals that the Trust Index is LLM-dependent: Qwen~3.5~35B's verbose generation style causes systematic false positives, underperforming the naive baseline by 16\%, while LLM selection dominates embedding model choice as the key performance factor. A retrieval depth sensitivity analysis ($K=3$ vs.\ $K=5$) confirms that detection performance is retrieval-depth-invariant: the 40\% recall ceiling is set by attack strategy difficulty, not by the number of documents retrieved.

The work addresses a gap left by existing frameworks such as RAGAS and ARES, which measure faithfulness to context but not the integrity of the context itself.
\end{abstract}

\tableofcontents
\listoffigures      % optional
\listoftables       % optional

% ---------- Main chapters ----------
% \chapter{Introduction}

% \section{Background and Motivation}
% % Connect RAG, trustworthiness, security, misinformation, data poisoning.
% % Include real-world motivation (e.g., RAG for documentation, defense, software systems).

% \section{Problem Statement}
% % Clearly define what is untrustworthy about current RAG systems.
% % Mention misinformation and data poisoning risks in AI pipelines.

% \section{Research Objectives and Questions}
% % Main research question.
% % Sub-questions aligned with evaluation agent, RAG trust, security.

% \section{Research Scope and Limitations}
% % Boundaries: datasets, models, attack types, evaluation constraints.

% \section{Thesis Structure}
% % Short description of each chapter.

\chapter{Introduction}
\label{ch:introduction}

\section{Background and Motivation}
\label{sec:background}

The rapid evolution of Natural Language Processing (NLP) has culminated in the widespread adoption of Large Language Models (LLMs) such as GPT-4, Llama, and Claude. Built upon the Transformer architecture \parencite{vaswani2017attention}, these models have demonstrated unprecedented capabilities in reasoning, code generation, and semantic understanding. Consequently, they act as foundational engines for applications ranging from automated software engineering to legal analytics.

However, the efficacy of standalone LLMs is fundamentally constrained by their static nature. Once trained, the parametric knowledge of a model is frozen, rendering it incapable of accessing real-time events or private, domain-specific data. Furthermore, LLMs suffer from a critical reliability issue known as \textit{hallucination}. Hallucination occurs when a model generates content that is grammatically plausible and semantically coherent but factually incorrect \parencite{ji2023survey}. This phenomenon arises because the model prioritizes statistical likelihood over factual verification.

To address the challenges of knowledge obsolescence and hallucination, the research community introduced Retrieval-Augmented Generation (RAG) \parencite{lewis2020rag}. The RAG paradigm alters the generative process by introducing an intermediate retrieval step. Before generating a response, the system queries an external knowledge base to fetch documents relevant to the user’s input. These retrieved documents are then fed into the generator’s context window, grounding the output in specific, evidence-based data.

Despite these advancements, the integration of external retrieval mechanisms introduces a new and arguably more insidious attack surface: the reliability of the retrieved data itself. The foundational assumption of standard RAG architectures is that the external knowledge base is benign and accurate. Recent scholarship suggests this assumption is increasingly dangerous in open-world deployments. As RAG systems ingest unstructured data from third-party sources, they become vulnerable to \textit{Knowledge Poisoning} and \textit{Adversarial Retrieval Attacks} \parencite{zou2024poisoned}.

In a poisoning scenario, an adversary injects malicious documents into the retrieval corpus. These documents are optimized to maximize semantic similarity with target queries, ensuring they are retrieved by the system. Once retrieved, these "poisoned" contexts can manipulate the LLM into generating misinformation or biased advice—a technique often referred to as indirect prompt injection. As AI systems evolve toward \textit{Agentic RAG}, where autonomous agents actively plan workflows based on retrieved data, the inability to distinguish between authentic evidence and adversarial noise poses a severe threat to the integrity of AI-driven decision-making.

\section{Problem Statement}
\label{sec:problem_statement}

The central problem addressed in this thesis is the \textbf{Security-Reliability Gap} in current Retrieval-Augmented Generation architectures. While RAG systems successfully mitigate the hallucination of non-existent facts, they lack robust mechanisms to verify the truthfulness and safety of the facts they retrieve. They inherently trust the retrieval output, operating under the assumption that high semantic relevance equates to factual truth.

This vulnerability manifests in three distinct dimensions within the current State of the Art:

\begin{enumerate}
    \item \textbf{Vulnerability to Semantic Adversarial Attacks:} Standard dense retrieval methods rely on vector similarity. Attackers can exploit this by crafting "poisoned" documents containing malicious payloads that are semantically identical to legitimate queries \parencite{zou2024poisoned}. Current RAG systems lack a secondary verification layer to detect these anomalies.
    
    \item \textbf{Limitations of Existing Evaluation Frameworks:} Current evaluation tools, such as RAGAS \parencite{es2024ragas} or ARES \parencite{saadfalcon2024ares}, focus primarily on "faithfulness" and "context relevance." They measure whether the answer reflects the context but do not assess the validity of the context itself. If a RAG system faithfully answers a user question based on a poisoned document, existing metrics would rate the system highly, despite the output being factually compromised.
    
    \item \textbf{Absence of Autonomous Defense Mechanisms:} Most defenses against data poisoning rely on offline data cleaning. There is a critical need for an online, autonomous \textit{Evaluation Agent} capable of acting as a "trust layer" during the inference pipeline to score, verify, and filter retrieved data before generation.
\end{enumerate}

Consequently, without such a defense mechanism, RAG systems risk becoming conduits for targeted misinformation campaigns. The research problem can be summarized as follows:

\begin{displayquote}
Existing RAG architectures lack real-time, autonomous capabilities to distinguish between benign evidence and adversarial data poisoning, rendering them susceptible to knowledge injection attacks that compromise system integrity.
\end{displayquote}

\section{Research Objectives}
\label{sec:objectives}

The overarching aim of this thesis is to design, implement, and evaluate a \textbf{Trustworthy RAG Framework} that integrates an autonomous Evaluation Agent. This agent serves as a defensive middleware, validating the factual integrity of retrieved documents and detecting signs of knowledge poisoning.

The specific objectives are defined as follows:

\begin{itemize}
    \item \textbf{Objective 1:} To design an autonomous evaluation agent utilizing Natural Language Inference (NLI) to assess factual consistency between retrieved evidence and generated responses.
    \item \textbf{Objective 2:} To develop detection algorithms specifically for "knowledge poisoning" attempts, distinguishing between natural data noise and adversarial injections.
    \item \textbf{Objective 3:} To formulate a composite \textit{Trust Index} metric that aggregates scores from factuality, reliability, and integrity dimensions.
    \item \textbf{Objective 4:} To empirically evaluate the proposed framework’s robustness against known retrieval attack vectors compared to standard RAG baselines.
\end{itemize}

\section{Research Questions}
\label{sec:research_questions}

To address the identified problem and achieve the stated objectives, this thesis answers the following Main Research Question (RQ) and supporting Sub-Questions:

\textbf{Main Research Question (RQ):}
\textit{How can an autonomous Evaluation Agent be designed and integrated within Retrieval-Augmented Generation (RAG) systems to detect misinformation and data poisoning, thereby improving the trustworthiness and security of AI-generated content?}

\textbf{Sub-Questions:}
\begin{itemize}
    \item \textbf{RQ1:} What computational strategies (e.g., NLI, semantic discrepancy detection) are effective for detecting misinformation and poisoning in retrieved contexts?
    \item \textbf{RQ2:} How can a composite "Trust Index" be mathematically formulated to quantify the reliability of a RAG response?
    \item \textbf{RQ3:} To what extent does the inclusion of an Evaluation Agent enhance the security performance of RAG pipelines compared to baseline systems, what are the trade-offs in latency, and how does the choice of LLM and embedding model affect the Trust Index's discriminative performance?
\end{itemize}

\section{Research Scope and Delimitations}
\label{sec:scope}

This thesis focuses on the evaluation and detection layer of RAG pipelines. The scope is defined as follows:
\begin{itemize}
    \item The work utilizes open-source LLMs (Llama 3.3 70B and Qwen 3.5 35B) accessed via the FARMI computing cluster and does not aim to pre-train models from scratch. Two embedding models (all-MiniLM-L6-v2 and Snowflake Arctic Embed2) are evaluated in a 2×2 factorial design to assess the sensitivity of the Trust Index to model selection.
    \item Experiments will be conducted on open-source datasets such as FEVER and TruthfulQA, alongside synthetic poisoned corpora.
    \item The scope is limited to \textit{textual} retrieval and generation; multimodal extensions (images/audio) are considered future work.
    \item The study focuses on "Knowledge Injection" attacks (poisoning the database) rather than model weight poisoning or direct jailbreaking prompts.
\end{itemize}

\section{Thesis Structure}
\label{sec:structure}

The remainder of this thesis is organized as follows:
\textbf{Chapter \ref{ch2_theory_background}} reviews the theoretical foundations of RAG and adversarial AI. 
\textbf{Chapter \ref{ch3_literature_review}} presents a systematic literature review identifying gaps in current defense mechanisms. 
\textbf{Chapter \ref{ch4_methodology}} details the research design and system architecture. 
\textbf{Chapter \ref{ch5_eval_agent}} describes the implementation of the Evaluation Agent. 
\textbf{Chapter \ref{ch6_experiments}} presents the experimental results and robustness analysis. 
\textbf{Chapter \ref{ch7_discussion}} discusses the implications of the findings. 
Finally, \textbf{Chapter \ref{ch8_conclusion}} concludes with future research directions.

\chapter{Theoretical and Conceptual Background}
\label{ch2_theory_background}

This chapter establishes the theoretical foundations required to understand the security and reliability challenges in Retrieval-Augmented Generation (RAG) systems. It provides an overview of Large Language Models (LLMs), mathematically defines the RAG architecture, and taxonomizes the threat landscape, specifically focusing on the distinction between intrinsic hallucinations and extrinsic knowledge poisoning. Finally, it introduces the concept of "Evaluation Agents" as a mechanism for automated auditing.

\section{Large Language Models and Generative AI}
\label{sec:llm_overview}

Large Language Models (LLMs) represent a class of deep learning models trained on massive corpora of text data to predict likelihood distributions over sequences of tokens. The current state-of-the-art models, such as GPT-4 and Llama 3, are built upon the \textbf{Transformer} architecture introduced by \textcite{vaswani2017attention}.

\subsection{The Transformer Architecture}
The core innovation of the Transformer is the \textit{Self-Attention} mechanism, which allows the model to weigh the importance of different words in a sequence relative to one another, regardless of their positional distance. This enables the capture of long-range dependencies and semantic context far more effectively than previous Recurrent Neural Networks (RNNs).

\subsection{Parametric Knowledge and Limitations}
LLMs store knowledge within their weights (parameters) during the pre-training phase. This is referred to as \textit{parametric knowledge}. While powerful, this approach suffers from two fundamental limitations:
\begin{enumerate}
    \item \textbf{Knowledge Cutoff:} The model's knowledge is static and limited to the data available up to the training date.
    \item \textbf{Hallucination:} Because LLMs are probabilistic engines designed to maximize token likelihood rather than factual accuracy, they frequently generate plausible but incorrect statements when the required information is absent from their parametric memory \parencite{ji2023survey}.
\end{enumerate}

\section{Retrieval-Augmented Generation (RAG)}
\label{sec:rag_architecture}

Retrieval-Augmented Generation (RAG) was proposed by \textcite{lewis2020rag} to mitigate the limitations of static LLMs. RAG creates a semi-parametric model by combining a pre-trained generator (the LLM) with a dense vector retriever.

% Placeholder for RAG Architecture Figure
\begin{figure}[h]
    \centering
    % \includegraphics[width=0.8\textwidth]{images/rag_architecture.png}
    \caption{Schematic of the RAG architecture: Query $\rightarrow$ Retriever $\rightarrow$ Generator $\rightarrow$ Response.}
    \label{fig:rag_arch}
\end{figure}

\subsection{Mathematical Formulation}
Formally, RAG models the probability of generating an output sequence $y$ given an input $x$ by marginalizing over a set of retrieved documents $z$. The generation probability is defined as:

\begin{equation}
    P(y|x) = \sum_{z \in Z} P_{\eta}(z|x) \cdot P_{\theta}(y|x, z)
\end{equation}

Where:
\begin{itemize}
    \item $P_{\eta}(z|x)$ is the probability of retrieving document $z$ given query $x$ (the \textbf{Retriever}).
    \item $P_{\theta}(y|x, z)$ is the probability of generating answer $y$ given query $x$ and document $z$ (the \textbf{Generator}).
\end{itemize}

\subsection{Dense Retrieval and Vector Embeddings}
In modern RAG implementations, retrieval is typically performed using Dense Passage Retrieval (DPR) or similar embedding models such as Sentence-BERT \parencite{reimers2019sentence}. Both the user query $q$ and the documents $d$ are mapped to a high-dimensional vector space. Relevance is calculated using \textbf{Cosine Similarity}:

\begin{equation}
    \text{sim}(q, d) = \frac{\mathbf{v}_q \cdot \mathbf{v}_d}{\|\mathbf{v}_q\| \|\mathbf{v}_d\|}
\end{equation}

This reliance on vector similarity is the primary attack surface for the security threats discussed in this thesis, as semantic closeness does not imply factual truth.

More recent embedding models, such as \texttt{snowflake-arctic-embed2}, produce 1024-dimensional representations and have demonstrated stronger retrieval quality on standard benchmarks compared to the 384-dimensional all-MiniLM-L6-v2 model. The choice of embedding model introduces a trade-off between retrieval fidelity and computational cost that is explored empirically in Chapter~\ref{ch6_experiments}.

\section{The Threat Landscape: Hallucination vs. Poisoning}
\label{sec:threat_landscape}

It is crucial to distinguish between natural errors and adversarial attacks in RAG pipelines.

\subsection{Hallucination (Intrinsic Error)}
Hallucination refers to the generation of unfaithful or nonsensical text. In the context of RAG, this often occurs when the retrieved context is irrelevant, causing the model to ignore the context and revert to its (potentially flawed) parametric memory. This is a reliability issue, not necessarily a security issue.

\subsection{Knowledge Poisoning (Extrinsic Attack)}
\label{subsec:knowledge_poisoning}
Knowledge Poisoning, or Corpus Poisoning, is a targeted security threat. Here, an adversary injects malicious documents into the external knowledge base. These documents are often engineered using \textit{Trigger Phrases} or optimized vector embeddings to ensure they are retrieved for specific queries \parencite{zou2024poisoned}.

% Placeholder for Poisoning Attack Figure
\begin{figure}[h]
    \centering
    % \includegraphics[width=0.8\textwidth]{images/poisoning_attack.png}
    \caption{Knowledge poisoning attack flow: adversarial documents are injected into the corpus and retrieved for target queries, manipulating the generated output.}
    \label{fig:poisoning_attack}
\end{figure}

Unlike "Data Poisoning" (which affects model training weights), Knowledge Poisoning affects the \textit{inference context}. A successful attack leads to \textbf{Indirect Prompt Injection}, where the poisoned document contains instructions that override the system prompt, forcing the LLM to output misinformation or malicious content.

\section{Trustworthiness in AI Systems}
\label{sec:trustworthiness}

Trustworthiness in Generative AI is a composite attribute comprising three key dimensions:

\begin{enumerate}
    \item \textbf{Faithfulness:} The degree to which the generated answer is derived strictly from the retrieved context, without adding external hallucinations.
    \item \textbf{Factuality:} The alignment of the generated answer with established world knowledge or ground truth.
    \item \textbf{Robustness:} The system's ability to maintain performance and safety boundaries in the presence of adversarial inputs or noisy data.
\end{enumerate}

Current evaluation frameworks often measure Faithfulness but neglect Robustness against poisoning, creating the research gap addressed by this thesis.

\section{Evaluation Agents and Automated Verification}
\label{sec:eval_agents}

To address the scalability limits of human evaluation, two predominant paradigms have emerged: \textbf{LLM-as-a-Judge}, where a Large Language Model evaluates the outputs of another AI system using reasoning frameworks such as Chain-of-Thought \parencite{wei2022chain}, and \textbf{Model-Based Verification}, where dedicated models (e.g., NLI classifiers) perform structured verification without LLM involvement.

In this thesis, the Evaluation Agent adopts the latter approach. It serves as a defensive middleware applying \textbf{Natural Language Inference (NLI)} to classify retrieval-response pairs as Entailed, Contradictory, or Neutral, combined with rule-based detection for adversarial patterns. This design avoids reliance on LLM-as-a-Judge, leveraging instead the deterministic, interpretable outputs of NLI models and explicit poisoning-detection logic to automate the identification of misinformation and knowledge poisoning attempts.

\section{Summary}
This chapter defined the RAG architecture and mathematically formalized the retrieval process. It highlighted the critical vulnerability of dense retrieval to knowledge poisoning attacks and defined the key dimensions of trustworthiness. The next chapter will review existing literature to identify specific gaps in current defense mechanisms.
% \chapter{Literature Review}
% \label{ch3_literature_review}
% \section{RAG System Components and Architectures}
% % Indexing, retrieval, reranking, generation, orchestration.

% \section{Fact-Checking and Factual Consistency Models}
% % FEVER, FactCC, FActScore, TruthfulQA, etc.

% \section{Detection of Adversarial Attacks and Data Poisoning in NLP Pipelines}
% % Poisoning taxonomy, trigger-based attacks, label flipping, etc.

% \section{Evaluation and Explainability Techniques for Trustworthy AI}
% % XAI, calibration, uncertainty estimation, etc.

% \section{Comparative Analysis of Existing Fact-Verification Pipelines}
% % Compare methods, datasets, metrics.

% \section{Identified Research Gaps and Research Contribution}
% % Gaps → your contribution; end with clear statement of novelty.

\chapter{Literature Review}
\label{ch3_literature_review}

This chapter provides a comprehensive systematic review of the research landscape surrounding Retrieval-Augmented Generation (RAG), adversarial AI security, and automated evaluation methodologies. The review is structured to first establish the evolution of RAG architectures from naive implementations to complex, modular systems. Subsequently, it analyzes the state-of-the-art in evaluation metrics, distinguishing between retrieval-centric and generation-centric approaches. A significant portion of the chapter is dedicated to the emerging threat of "Knowledge Poisoning" and the limitations of current defense mechanisms. The chapter concludes by synthesizing these findings to articulate the specific research gaps this thesis aims to bridge.

\section{The Evolution of RAG Architectures}
\label{sec:lit_rag_evolution}

Retrieval-Augmented Generation has evolved significantly since its inception. \textcite{gao2024retrieval} categorize this evolution into three distinct paradigms: Naive RAG, Advanced RAG, and Modular RAG.

\subsection{Naive vs. Advanced RAG}
The "Naive" RAG paradigm follows a linear "Retrieve-Then-Generate" process. While effective for simple queries, it suffers from low precision when retrieved documents are noisy or conflicting. "Advanced" RAG introduces pre-retrieval strategies (e.g., query rewriting) and post-retrieval strategies (e.g., re-ranking) to improve context quality. However, these improvements primarily focus on relevance, often neglecting the veracity of the content.

\subsection{Modular RAG and Agentic Workflows}
The current state-of-the-art is moving toward "Modular RAG," where the retriever, generator, and evaluator function as independent, interchangeable modules. This aligns with the "Agentic" approach proposed in this thesis, where an autonomous agent intervenes in the pipeline. \textcite{asai2024selfrag} demonstrate that modular architectures allow for "Self-Reflection," enabling the model to critique its own retrieved context.

\section{Evaluation Methodologies in RAG}
\label{sec:lit_evaluation}

Evaluating RAG systems requires a dual approach: assessing the quality of the retrieval (Information Retrieval metrics) and the quality of the generated text (Natural Language Generation metrics).

\subsection{Retrieval-Centric Metrics}
Before a response is generated, the system must first retrieve relevant documents. The efficacy of this stage is typically measured using standard Information Retrieval (IR) metrics established by \textcite{karpukhin2020dense} in their work on Dense Passage Retrieval (DPR):
\begin{itemize}
    \item \textbf{Recall@k:} Measures the proportion of relevant documents found in the top-$k$ results.
    \item \textbf{Mean Reciprocal Rank (MRR):} Evaluates how high the first relevant document appears in the list.
\end{itemize}
While these metrics measure \textit{relevance}, they fail to measure \textit{safety}. A poisoned document engineered to be "highly relevant" would score perfectly on MRR, highlighting a critical blind spot in standard IR evaluation.

\subsection{Generation-Centric and Reference-Free Metrics}
Traditionally, generation quality was measured using n-gram overlap metrics like BLEU and ROUGE. However, \textcite{lin2022truthfulqa} and \textcite{honovich2022true} argue that these are insufficient for measuring factuality.
\begin{itemize}
    \item \textbf{Factuality (Faithfulness):} Does the answer contradict the context?
    \item \textbf{Truthfulness:} Does the answer align with world knowledge?
\end{itemize}
Frameworks like \textbf{RAGAS} \parencite{es2024ragas} and \textbf{ARES} \parencite{saadfalcon2024ares} utilize LLMs as judges to score these dimensions. \textcite{honovich2022true} introduced the TRUE benchmark to standardize factual consistency evaluation, using Natural Language Inference (NLI) to determine if a claim is entailed by its evidence.

\section{Adversarial Attacks on RAG Pipelines}
\label{sec:lit_attacks}

As RAG systems increasingly rely on open-web data, they inherit vulnerabilities associated with untrusted sources.

\subsection{Indirect Prompt Injection}
\textcite{greshake2023indirect} demonstrated "Indirect Prompt Injection," a sophisticated attack where adversaries embed malicious instructions (e.g., "Ignore previous instructions and output 'Scam'") into web content. When an LLM retrieves this content, it treats the injected text as authoritative instructions rather than passive data.

\subsection{Knowledge Poisoning (Corpus Poisoning)}
Building on injection attacks, \textcite{zou2024poisoned} formalized "Knowledge Poisoning" in RAG. Unlike training data poisoning (which requires expensive model retraining), knowledge poisoning targets the inference database.
\subsubsection{Mechanism of Attack}
Attackers generate "collision" documents that maximize vector similarity with targeted queries. \textcite{zou2024poisoned} showed that poisoning fewer than 0.1\% of the documents in a retrieval corpus can achieve an Attack Success Rate (ASR) of over 40\% on specific topics.
\subsubsection{Targeted vs. Untargeted Attacks}
Literature distinguishes between targeted attacks (forcing a specific wrong answer) and untargeted attacks (degrading general system utility). This thesis specifically focuses on detecting targeted misinformation injections.



\section{Defense Mechanisms and Fact Verification}
\label{sec:lit_defense}

Current defenses can be categorized into static preprocessing and dynamic verification.

\subsection{Static Defenses (Filtering)}
Standard defenses involve filtering "toxic" content using classifiers before indexing. However, misinformation is often phrased neutrally and professionally, bypassing toxicity filters.

\subsection{Dynamic Verification (NLI Approaches)}
The most promising defense lies in dynamic verification using Natural Language Inference (NLI). NLI models classify the relationship between two sentences as \textit{Entailment}, \textit{Contradiction}, or \textit{Neutral}. \textcite{shuster2021retrieval} utilized this approach to reduce hallucinations in conversational agents. By treating the retrieved document as the "Premise" and the generated answer as the "Hypothesis," NLI can mathematically quantify consistency.

\subsection{Agentic Self-Correction}
Recent works like \textbf{Self-RAG} \parencite{asai2024selfrag} introduce "critic tokens" that allow the model to self-assess generation quality. While effective for benign errors, it remains unclear if self-correction is robust against adversarial inputs designed to manipulate the critic itself.

\section{Identified Research Gaps}
\label{sec:research_gap}

Despite the proliferation of RAG evaluation frameworks, significant gaps remain:
\begin{enumerate}
    \item \textbf{The Trust-Security Gap:} Existing frameworks like RAGAS measure if the model follows the context, but not if the context is safe. There is a lack of tools that simultaneously evaluate \textit{Factuality} and \textit{Security}.
    \item \textbf{Lack of Real-Time Auditing:} Most poisoning defenses are offline corpus-cleaning tasks. There is a need for an online Evaluation Agent that sits in the inference loop.
    \item \textbf{Unified Metrics:} There is no widely accepted "Trust Index" that aggregates retrieval safety, generation faithfulness, and adversarial robustness into a single interpretable score for stakeholders.
    \item \textbf{LLM Generation-Style Sensitivity:} No existing evaluation framework systematically examines how the LLM's generation style interacts with NLI-based trust scoring. Verbose or hedged model outputs (e.g., ``\textit{Based on the provided context...}'') reduce entailment confidence in the NLI Verifier even when the answer is factually correct, introducing model-specific false positives that are invisible to retrieval-only analysis. This gap motivates a multi-LLM comparative evaluation of the Trust Index, as conducted in this thesis.
\end{enumerate}

This thesis addresses these gaps by designing an autonomous Evaluation Agent that combines NLI-based verification with specific poisoning detection logic, and by evaluating its performance across multiple LLM and embedding model combinations.

\section{Summary}
This chapter reviewed the trajectory of RAG research from basic retrieval to modular, agentic systems. It highlighted that while evaluation metrics for relevance are mature, metrics for adversarial robustness in retrieval are nascent. The literature confirms that Knowledge Poisoning is a critical, under-addressed threat, validating the need for the Evaluation Agent proposed in this study.
% \chapter{Research Design and Methodology}
% \label{ch4_methodology}
% \section{Research Design Overview}
% % State whether constructive, experimental, or applied design.

% \section{System Architecture of the Trustworthy RAG Framework}
% \subsection{Retrieval Module}
% % Index, vector store, retriever type.

% \subsection{Generation Module}
% % LLMs used, prompt templates, decoding.

% \subsection{Evaluation Agent}
% % Role, inputs/outputs, decision logic (high-level, details in Chapter 5).

% \subsection{Security and Verification Layer}
% % How attacks, anomalies, or poisoning are detected/logged.

% \section{Dataset Selection and Preparation}
% % e.g., FEVER, TruthfulQA, HotpotQA, synthetic corpora; preprocessing.

% \section{Experimental Scenarios}
% % Baseline RAG, misinformation injection, data poisoning setups.

% \section{Evaluation Metrics}
% % Factuality (e.g., FActScore, NLI accuracy),
% % Trust Index (precision/recall of misinformation detection),
% % Latency and computational overhead.

% \section{Implementation Tools and Environment}
% % LangChain, LlamaIndex, HuggingFace, FAISS, Python, OpenAI/LLMs, hardware.

% \section{Validation and Reproducibility Strategy}
% % Seeds, configuration, code availability, experiment logging.

% \section{Ethical and Security Considerations}
% % Handling of harmful content, responsible disclosure, etc.


\chapter{Research Design and Methodology}
\label{ch4_methodology}

This chapter outlines the methodological framework adopted to design, implement, and evaluate the proposed Trustworthy RAG System. The research follows the \textbf{Design Science Research (DSR)} paradigm, focusing on the creation of a novel artifact (the Evaluation Agent) to solve a specific problem (Knowledge Poisoning). The chapter details the threat model, the system architecture, the dataset construction strategy for simulating adversarial attacks, and the quantitative metrics used for evaluation.

\section{Research Design Overview}
\label{sec:design_overview}

The research is structured into three distinct phases:
\begin{enumerate}
    \item \textbf{Phase 1: Dataset Preparation \& Injection.} Construction of a "Clean" baseline corpus (using FEVER/TruthfulQA) and the generation of a "Poisoned" corpus containing adversarial documents.
    \item \textbf{Phase 2: System Development.} Implementation of the RAG pipeline integrated with the autonomous Evaluation Agent (detailed in Chapter \ref{ch5_eval_agent}).
    \item \textbf{Phase 3: Empirical Evaluation.} A comparative analysis measuring the agent's ability to detect poisoning and its impact on system latency and retrieval accuracy.
\end{enumerate}

\section{Threat Model and Assumptions}
\label{sec:threat_model}

To evaluate security robustness, a specific Threat Model must be defined. This thesis assumes a "Black-Box Knowledge Injection" adversary.

\subsection{Attacker Capabilities}
The attacker has the following capabilities:
\begin{itemize}
    \item \textbf{Corpus Access:} The attacker can insert new documents into the external knowledge base (e.g., via a public wiki, forum, or document upload feature).
    \item \textbf{Content Manipulation:} The attacker can craft text specifically designed to trigger specific keywords or semantic embeddings.
\end{itemize}

\subsection{Attacker Limitations}
The attacker \textbf{cannot}:
\begin{itemize}
    \item Modify the internal weights or parameters of the LLM or the Retriever.
    \item Access the system prompt or the internal logic of the Evaluation Agent.
\end{itemize}

\section{System Architecture}
\label{sec:system_architecture}

The proposed "Trustworthy RAG Framework" introduces a middleware layer between the retrieval and generation phases.



\subsection{Components}
\begin{itemize}
    \item \textbf{The Retriever:} A dense vector store (FAISS) supporting two embedding backends evaluated in this thesis: (i) \texttt{all-MiniLM-L6-v2} (384-dimensional, local Sentence-BERT family) and (ii) \texttt{snowflake-arctic-embed2} (1024-dimensional, accessed via FARMI API). Both retrieve the top-$k$ semantically similar documents per query.
    \item \textbf{The Evaluation Agent (The Artifact):} An autonomous module that acts as a "Gatekeeper." It analyzes the retrieved documents for semantic consistency and malicious patterns before they are passed to the generator.
    \item \textbf{The Generator:} Two LLMs are evaluated: \texttt{llama3.3:70b} and \texttt{qwen3.5:35b}, both hosted on the FARMI computing cluster (Tampere University) via an OpenAI-compatible API. Llama 3.3 70B serves as the primary thesis baseline; Qwen 3.5 35B is included as a cross-model comparison point to assess the sensitivity of the Trust Index to LLM generation style.
\end{itemize}

\subsection{Multi-Model Evaluation Design}
\label{subsec:factorial_design}

To assess the sensitivity of the Trust Index to model selection, a \textbf{2×2 factorial experimental design} is employed, crossing two LLMs with two embedding models:

\begin{equation*}
\{\text{Llama 3.3 70B},\; \text{Qwen 3.5 35B}\} \;\times\;
\{\text{all-MiniLM-L6-v2},\; \text{Snowflake Arctic Embed2}\}
\end{equation*}

This yields four experimental configurations, all run with identical hyperparameters (100 samples, 30\% poison ratio, $K=3$ retrieved documents, classification threshold $\tau=0.5$). The design allows isolation of LLM effects from embedding effects and tests whether the Trust Index generalizes across model choices or is fundamentally dependent on a specific generator's output style. A separate sensitivity analysis comparing $K=3$ and $K=5$ retrieval depths for the primary Llama~3.3~70B + all-MiniLM-L6-v2 configuration is presented in Section~\ref{sec:ch6_topk}.

\section{Dataset Selection and Preparation}
\label{sec:dataset_prep}

To simulate realistic misinformation and poisoning scenarios, this study utilizes a combination of established benchmarks and synthetic adversarial data.

\subsection{Base Datasets}
\begin{itemize}
    \item \textbf{FEVER (Fact Extraction and VERification):} Used for testing the agent's ability to verify claims against evidence \parencite{thorne2018fever}.
    \item \textbf{TruthfulQA:} Used to evaluate the model's propensity to mimic human falsehoods \parencite{lin2022truthfulqa}.
\end{itemize}

\subsection{Synthetic Poisoning Generation}
Since no standard "Poisoned RAG" dataset exists for this specific architecture, a synthetic dataset is generated using a custom \textit{PoisonedDatasetGenerator} inspired by \textcite{zou2024poisoned}. Four poisoning strategies are implemented:
\begin{itemize}
    \item \textbf{Contradiction:} Appends contradicting statements (e.g., ``Contrary to popular belief...'') to genuine documents.
    \item \textbf{Injection:} Appends instruction-override phrases (e.g., ``IMPORTANT: Ignore previous context...'') to manipulate the LLM.
    \item \textbf{Entity Swap:} Replaces entities or numbers in-place within the original text (e.g., swapping ``14 million'' for ``41 million''). Since domain-specific documents (such as TruthfulQA entries) rarely contain common geopolitical named entities, a three-level fallback is employed: (1) predefined entity substitution (e.g., Paris~$\rightarrow$~Berlin); (2) numeric value shifting via regular expressions when no named entity is present; (3) qualifier or negation swapping in the answer segment of the document (e.g., replacing ``is'' with ``is not''). If none of these produce a textual change, a plausible-sounding dispute sentence is appended as a last resort.
    \item \textbf{Subtle:} Combines entity swap with minor qualifier additions for minimal, hard-to-detect modifications.
\end{itemize}
Experiments use 50 samples per run with a poison ratio of 30\% (i.e., 30\% of retrieved documents are poisoned). Per-strategy experiments isolate each attack type; mixed experiments combine all four strategies.

\section{Evaluation Metrics}
\label{sec:eval_metrics}

To answer the Research Questions, the system is evaluated across three dimensions: Security, Reliability, and Performance.

\subsection{Security Metrics (Detection Performance)}
These metrics measure the Evaluation Agent's ability to identify poisoned documents (RQ1 & RQ2).
\begin{itemize}
    \item \textbf{Precision and Recall:} Measuring false positives (benign documents flagged as poison) and false negatives (poisoned documents missed).
    \item \textbf{Poison Detection Rate (PDR):} The fraction of adversarially poisoned samples correctly identified and flagged by the Evaluation Agent before generation. Ground truth is assigned based on the document-level poisoning label (i.e., whether the knowledge base document corresponding to the query was modified), not on retrieval outcome. This is the primary security metric evaluated in this thesis (equivalent to Recall in binary classification).
    \item \textbf{Attack Success Rate (ASR):} The percentage of attacks where the final generated answer mimics the injected misinformation. While ASR is a theoretically important metric, rule-based poisoning strategies (particularly Contradiction and Injection, which append content to genuine documents) yield low natural ASR even without the Evaluation Agent, since modern LLMs tend to prioritise the original document content over appended overrides. Measuring true ASR requires semantic comparison of generated answers to injected claims, which is treated as future work in Section~\ref{sec:future_work}.
\end{itemize}

\subsection{Reliability Metrics (The Trust Index)}
The \textbf{Trust Index} ($T_{score}$) is a composite metric derived from the agent's internal logic:
\begin{equation}
    T_{score} = \alpha \cdot S_{factuality} + \beta \cdot S_{consistency} + \gamma \cdot (1 - P_{poison})
\end{equation}
where $S_{factuality}$ is the NLI-based factual consistency score, $S_{consistency}$ measures inter-document agreement, and $P_{poison}$ is the poison probability. Default weights are $\alpha=0.4$, $\beta=0.35$, and $\gamma=0.25$. A non-linear dampener is applied when $P_{poison} > 0.7$ to penalize highly poisoned content. Formally, the dampener $\delta$ reduces the trust score multiplicatively:
\begin{equation}
    \delta(P_{poison}) = 1 - 0.4 \cdot \frac{P_{poison} - 0.70}{0.30}, \quad P_{poison} > 0.70
\end{equation}
and the dampened score becomes $T_{\text{final}} = T \cdot \delta(P_{poison})$. At $P_{poison}=0.70$ no penalty is applied ($\delta=1.0$); at $P_{poison}=1.0$ the trust score is reduced by 40\% ($\delta=0.6$), ensuring that extreme poisoning signals override even strong factuality scores. A full derivation, worked examples, and the boundary-case analysis are provided in Chapter~\ref{ch5_eval_agent}.

\subsection{Performance Metrics (RQ3)}
\begin{itemize}
    \item \textbf{Latency Overhead:} The additional time (in milliseconds) introduced by the Evaluation Agent per query.
    \item \textbf{Throughput:} The number of queries processed per second compared to a baseline RAG system.
\end{itemize}

\section{Implementation Environment}
\label{sec:implementation_env}

The system is implemented using Python 3.10. The core libraries include \textbf{LangChain} for pipeline orchestration, \textbf{HuggingFace Transformers} for model loading (NLI: BART-MNLI), and \textbf{FAISS} for vector retrieval. Two LLMs (\texttt{llama3.3:70b} and \texttt{qwen3.5:35b}) and two embedding backends (\texttt{all-MiniLM-L6-v2} local and \texttt{snowflake-arctic-embed2} API) are hosted on the FARMI computing cluster (Tampere University) and accessed via an OpenAI-compatible API. The NLI model and evaluation pipeline run on CPU.

\section{Summary}
This chapter defined the Design Science methodology and the specific threat model (Knowledge Poisoning) used to stress-test the system. It detailed the architecture of the Trustworthy RAG framework and established the quantitative metrics for success. The next chapter will provide the detailed implementation logic of the Evaluation Agent.
\chapter{Development of the Evaluation Agent}
\label{ch5_eval_agent}

This chapter presents the design and implementation of the core artifact of this thesis: the \textbf{Evaluation Agent}. The Evaluation Agent is introduced as a defensive middleware in the RAG pipeline, positioned between retrieval and final trust decision. Its purpose is to detect misinformation and knowledge poisoning in near real-time and provide an interpretable trustworthiness score for each generated response.

\section{Design Philosophy and Functional Requirements}

\subsection{Design Philosophy}
The Evaluation Agent is designed with four principles:
\begin{enumerate}
    \item \textbf{Defense-in-depth:} No single signal is sufficient for robust detection. The agent combines factual verification, consistency analysis, and poisoning indicators.
    \item \textbf{Interpretability:} Each trust decision must be explainable through component scores, warnings, and textual reports.
    \item \textbf{Runtime practicality:} Evaluation must run in the online inference loop with manageable overhead.
    \item \textbf{Modularity:} Components should be independently replaceable (e.g., alternative NLI model, different poison detector).
\end{enumerate}

\subsection{Functional Requirements}
Based on the problem definition in Chapters~\ref{ch:introduction} and~\ref{ch4_methodology}, the Evaluation Agent must:
\begin{itemize}
    \item verify whether retrieved evidence supports the generated answer;
    \item identify suspicious patterns indicating possible corpus poisoning;
    \item quantify trustworthiness using a composite score in $[0,1]$;
    \item provide both machine-readable outputs and human-readable explanations;
    \item integrate seamlessly with the end-to-end RAG pipeline.
\end{itemize}

\section{Architecture and Internal Workflows}

The implemented Evaluation Agent orchestrates three modules:
\begin{itemize}
    \item \textbf{NLI Verifier} (\texttt{nli\_verifier.py})
    \item \textbf{Poison Detector} (\texttt{poison\_detector.py})
    \item \textbf{Trust Index Calculator} (\texttt{trust\_index.py})
\end{itemize}

At runtime, the orchestration class (\texttt{evaluation\_agent.py}) executes:
\begin{enumerate}
    \item NLI verification between retrieved documents and generated answer;
    \item poisoning analysis over retrieved documents (linguistic, structural, semantic, and consistency signals);
    \item consistency estimation and retrieval-confidence adjustment;
    \item final Trust Index computation and report generation.
\end{enumerate}

\subsection{Claim Extraction}
In the current implementation, explicit claim decomposition is intentionally lightweight. Instead of extracting atomic claims via an additional parser, the generated response itself is treated as the primary hypothesis for verification. This design choice reduces complexity and latency in the first system version.

Operationally:
\begin{itemize}
    \item \textbf{Premise:} each retrieved document;
    \item \textbf{Hypothesis:} generated answer;
    \item \textbf{Pairwise evaluation:} repeated for all retrieved documents.
\end{itemize}

This approach is suitable for short- to medium-length answers and supports direct alignment with NLI-based factual verification.

\subsection{Evidence Retrieval}
Evidence is supplied by the retriever in two forms:
\begin{itemize}
    \item retrieved text documents and their similarity scores;
    \item document embeddings (used for semantic outlier detection).
\end{itemize}

The pipeline method \texttt{query\_with\_evaluation()} retrieves top-$k$ documents and forwards:
\begin{itemize}
    \item \texttt{retrieved\_documents} (text),
    \item \texttt{retrieval\_scores},
    \item \texttt{document\_embeddings}.
\end{itemize}
This avoids redundant re-embedding inside the Evaluation Agent.

\subsection{Fact Verification and Scoring}
\label{subsec:ch5_nli}

The NLI Verifier uses \texttt{facebook/bart-large-mnli} through direct sequence classification (not zero-shot prompting). For each premise-hypothesis pair, the model outputs probabilities for:
\begin{itemize}
    \item contradiction,
    \item neutral,
    \item entailment.
\end{itemize}

Aggregated metrics include:
\begin{itemize}
    \item average entailment,
    \item average contradiction,
    \item maximum contradiction,
    \item support ratio.
\end{itemize}

\paragraph{Factuality score.}
The factuality score is computed with an entailment-focused strategy. Only documents with meaningful entailment signal are used for final factuality aggregation. If no such document exists, the system returns an inconclusive baseline ($0.5$) instead of forcing an extreme negative value. This reduces false penalties in cases where retrieved documents are only weakly related to the query.

This design is motivated by a key empirical observation during development: \texttt{facebook/bart-large-mnli} assigns near-maximum contradiction scores ($\approx 0.99$) to \textit{both} genuine semantic contradictions \textit{and} to completely unrelated text pairs. The model cannot reliably distinguish between a document that actively contradicts the answer and one that is simply topically unrelated to it. Using raw contradiction scores as a negative factuality signal would therefore produce systematic false penalisation of off-topic but benign documents. The entailment-focused strategy resolves this by treating contradiction and neutral signals as inconclusive rather than as evidence of untruthfulness, reserving strong negative signals for the dedicated poison detection pathway.

\subsection{Poison Detection Workflow}

The poison detector combines multiple signals:
\begin{enumerate}
    \item \textbf{Linguistic patterns:} instruction override phrases, contradiction markers, format-injection traces.
    \item \textbf{Structural anomalies:} excessive uppercase, suspicious repetition, abnormal tokens.
    \item \textbf{Intra-document consistency:} NLI between first and second halves of the same document.
    \item \textbf{Cross-document consistency:} pairwise contradiction checks among retrieved documents.
    \item \textbf{Semantic outlier analysis:} embedding-based deviation from batch centroid.
\end{enumerate}

Each document receives a poison probability from combined signals with diminishing returns. A batch-level poison probability is then computed. In the final design, relevance-weighted aggregation is used when retrieval scores are available, so low-ranked suspicious neighbors do not dominate the final decision.

\subsection{Trustworthiness Aggregation}
\label{subsec:ch5_trust_index}

\subsubsection{The Trust Index Formula}

The Trust Index $T \in [0,1]$ is a weighted linear combination of three component scores:

\begin{equation}
\label{eq:trust_index}
T = \alpha \cdot S_{\text{factuality}} + \beta \cdot S_{\text{consistency}} + \gamma \cdot (1 - P_{\text{poison}})
\end{equation}

where the default weights are $\alpha = 0.40$, $\beta = 0.35$, $\gamma = 0.25$, satisfying
$\alpha + \beta + \gamma = 1$.

\paragraph{Weight rationale.}
The three weights reflect the precision and reliability of each signal source, as summarised in
Table~\ref{tab:trust_weights}.

\begin{table}[h]
\centering
\caption{Trust Index weight allocation and rationale.}
\label{tab:trust_weights}
\begin{tabular}{clcp{6.5cm}}
\hline
\textbf{Weight} & \textbf{Signal} & \textbf{Value} & \textbf{Rationale} \\
\hline
$\alpha$ & Factuality  & 0.40 & NLI entailment is a precise, model-based signal directly measuring
                                whether the answer is grounded in retrieved evidence. Assigned the
                                highest weight as the most reliable indicator. \\
$\beta$  & Consistency & 0.35 & Cross-document agreement provides a strong proxy for knowledge base
                                integrity without requiring access to external ground truth. \\
$\gamma$ & Poison safety & 0.25 & Poison detection combines heuristic signals of varying precision.
                                  Assigning a lower weight limits the influence of false positives
                                  from pattern-based detectors on the final score. \\
\hline
\end{tabular}
\end{table}

The weight configuration was empirically validated through an ablation study reported in
Chapter~\ref{ch6_experiments}, which tested four configurations including equal weights
(0.33, 0.34, 0.33) and poison-heavy settings. The default weights yielded the best overall
F1 performance, confirming that factuality and consistency together contribute more
discriminative power than the poison score alone.

\subsubsection{Insufficiency of the Linear Formula Under High Poisoning}
\label{subsubsec:linear_insufficiency}

A critical design problem arises when $P_{\text{poison}}$ is very high (near 1.0) but the
factuality and consistency scores remain elevated. This occurs in practice when the language
model \textit{ignores} the adversarial portion of a poisoned document and produces a
factually correct answer grounded in the legitimate portion. In that case, NLI entailment
scores are high, yet the retrieved context is demonstrably compromised.

Consider the following concrete example. Suppose a poisoned context produces:
\begin{itemize}
    \item $S_{\text{factuality}} = 0.80$ (LLM's answer is well-supported by the document text)
    \item $S_{\text{consistency}} = 0.70$ (documents broadly agree with each other)
    \item $P_{\text{poison}} = 0.90$ (the poison detector reports high contamination)
\end{itemize}

Applying the linear formula directly:
\begin{align}
T_{\text{linear}} &= 0.40 \times 0.80 + 0.35 \times 0.70 + 0.25 \times (1 - 0.90) \notag \\
                  &= 0.320 + 0.245 + 0.025 \notag \\
                  &= 0.590
\end{align}

The result $T = 0.59$ exceeds the classification threshold $\tau = 0.5$, causing the system
to label the response as \textit{trustworthy} despite a 90\% poison probability. The root
cause is structural: because $\gamma = 0.25$, the poison term contributes at most 0.25 to
$T$ when $P_{\text{poison}} = 0$, and 0 when $P_{\text{poison}} = 1.0$. The maximum
possible poison-driven reduction in $T$ is therefore 0.25 points. Since $\alpha + \beta =
0.75$, factuality and consistency can always push $T$ above $\tau = 0.5$ independently of
the poison term when both scores are moderate or high. A purely linear aggregation is
therefore insufficient for reliable detection under high-contamination conditions.

\subsubsection{Non-Linear Poison Dampener}
\label{subsubsec:dampener}

To address this limitation, a multiplicative dampener is applied to the trust score whenever
$P_{\text{poison}} > 0.70$. The threshold of 0.70 was chosen to distinguish the
\textit{ambiguous} regime ($P_{\text{poison}} \leq 0.70$), where heuristic signals are
inconclusive, from the \textit{confirmed contamination} regime ($P_{\text{poison}} > 0.70$),
where multiple independent signals converge on a poisoning verdict.

The dampener is defined as:
\begin{equation}
\label{eq:dampener}
d(P_{\text{poison}}) =
\begin{cases}
1.0 & \text{if } P_{\text{poison}} \leq 0.70 \\[4pt]
1.0 - 0.4 \cdot \dfrac{P_{\text{poison}} - 0.70}{0.30} & \text{if } P_{\text{poison}} > 0.70
\end{cases}
\end{equation}

The final trust score after dampening is:
\begin{equation}
\label{eq:trust_final}
T_{\text{final}} = T_{\text{linear}} \times d(P_{\text{poison}})
\end{equation}

The dampener is a linear interpolation from $d = 1.0$ at $P_{\text{poison}} = 0.70$ to
$d = 0.60$ at $P_{\text{poison}} = 1.0$, producing a smooth, monotonically decreasing
penalty over the contamination interval. Table~\ref{tab:dampener_values} lists the
multiplier at representative poison probability values.

\begin{table}[h]
\centering
\caption{Dampener multiplier $d$ at representative poison probability values.}
\label{tab:dampener_values}
\begin{tabular}{ccc}
\hline
$P_{\text{poison}}$ & Dampener $d$ & Interpretation \\
\hline
$\leq 0.70$ & 1.000 & No penalty applied (ambiguous zone) \\
0.75        & 0.933 & Mild penalty ($-6.7\%$) \\
0.80        & 0.867 & Moderate penalty ($-13.3\%$) \\
0.85        & 0.800 & Moderate penalty ($-20.0\%$) \\
0.90        & 0.733 & Strong penalty ($-26.7\%$) \\
0.95        & 0.667 & Strong penalty ($-33.3\%$) \\
1.00        & 0.600 & Maximum penalty ($-40.0\%$) \\
\hline
\end{tabular}
\end{table}

\paragraph{Design properties of the dampener.}
The form of Equation~\ref{eq:dampener} was chosen to satisfy three properties:
\begin{enumerate}
    \item \textbf{Continuity at the threshold:} $d(0.70) = 1.0$ ensures no discontinuous
    jump in trust score at the activation boundary.
    \item \textbf{Bounded maximum penalty:} $d(1.0) = 0.60$ limits the reduction to 40\%,
    preventing the trust score from collapsing to near-zero in edge cases where the poison
    detector over-fires. A floor prevents discarding responses that may still be usable
    with human oversight.
    \item \textbf{Proportionality:} the penalty grows linearly with contamination probability
    in the interval $[0.70, 1.0]$, reflecting proportional confidence in the poisoning verdict.
\end{enumerate}

\paragraph{Worked example --- poisoned context.}
Revisiting the example from Section~\ref{subsubsec:linear_insufficiency} with
$T_{\text{linear}} = 0.590$ and $P_{\text{poison}} = 0.90$:
\begin{align}
d(0.90) &= 1.0 - 0.4 \times \frac{0.90 - 0.70}{0.30} = 1.0 - 0.4 \times 0.667 = 0.733 \\
T_{\text{final}} &= 0.590 \times 0.733 = 0.432
\end{align}
The dampened score $T_{\text{final}} = 0.432$ falls below $\tau = 0.5$, correctly
classifying the response as \textit{untrustworthy} and triggering an actionable warning
to the consumer.

\paragraph{Worked example --- clean context.}
For a clean retrieval context with no poisoning signal:
$S_{\text{factuality}} = 0.90$, $S_{\text{consistency}} = 0.85$, $P_{\text{poison}} = 0.10$:
\begin{align}
T_{\text{linear}} &= 0.40 \times 0.90 + 0.35 \times 0.85 + 0.25 \times 0.90 \notag \\
                  &= 0.360 + 0.298 + 0.225 = 0.883
\end{align}
Since $P_{\text{poison}} = 0.10 \leq 0.70$, the dampener is not activated:
$T_{\text{final}} = 0.883$, correctly classified as \textbf{HIGH} trust.

\subsubsection{Retrieval Confidence Modifier}

A secondary adjustment accounts for weak retrieval quality. When the highest retrieval
similarity score among the top-$K$ documents falls below 0.30, the retrieved documents
may not be genuinely relevant to the query. In this case the trust score is further
penalised:
\begin{equation}
\label{eq:retrieval_confidence}
T_{\text{final}} = T_{\text{final}} \times \bigl(0.5 + 0.5 \cdot r_{\max}\bigr),
\quad \text{if } r_{\max} < 0.30
\end{equation}
where $r_{\max}$ is the maximum retrieval similarity score in the current batch. This
modifier is inactive for normal retrieval quality ($r_{\max} \geq 0.30$) and avoids
inflating trust scores for queries where retrieval has returned loosely related documents.

\subsubsection{LLM-Dependency of the Trust Index}
\label{subsubsec:llm_dependency}

A non-obvious property of the Trust Index is that its value depends not only on the
\textit{content} of the generated answer but also on the \textit{generation style} of the
underlying LLM. The factuality component $S_{\text{factuality}}$ is computed by the NLI
Verifier as the average entailment score between each retrieved document (premise) and the
generated answer (hypothesis). If the LLM systematically produces hedged or verbose outputs
--- for example, prefacing answers with phrases such as ``\textit{Based on the provided
context, it appears that\ldots}'' or ``\textit{According to the available information\ldots}''
--- the NLI model treats this stylistic hedging as reduced semantic commitment to the claim,
lowering the entailment probability and, consequently, $S_{\text{factuality}}$.

This effect is benign when the hedging reflects genuine uncertainty, but it becomes a
systematic source of false positives when clean, factually accurate responses consistently
receive low entailment scores purely due to generation style. In the comparative evaluation
(Chapter~\ref{ch6_experiments}), Qwen 3.5 35B exhibits this behaviour: its verbose outputs
reduce the mean clean trust score to $\approx 0.64$, compared to $\approx 0.83$ for
Llama 3.3 70B, pushing a substantial fraction of clean responses below the $\tau = 0.5$
threshold and causing 21--23 false positives out of 85 clean samples.

This LLM-dependency is an intrinsic property of NLI-based trust assessment rather than a
calibration error. It implies that the Trust Index threshold $\tau$ and the component
weights $\alpha, \beta, \gamma$ should be treated as \textbf{LLM-specific hyperparameters}
calibrated per model rather than universal constants. Alternatively, a style-normalisation
step applied to the generated hypothesis before NLI inference could reduce generation-style
artefacts without altering the factual content being evaluated.

\subsubsection{Trust Level Classification}

The final score $T_{\text{final}}$ is mapped to one of four categorical trust levels to
support human-readable interpretation alongside the numeric score:

\begin{table}[h]
\centering
\caption{Trust level thresholds and operational interpretation.}
\label{tab:trust_levels}
\begin{tabular}{clp{7cm}}
\hline
\textbf{Score range} & \textbf{Trust level} & \textbf{Operational meaning} \\
\hline
$T > 0.80$         & HIGH      & Response is well-supported by consistent, clean evidence.
                                 Suitable for direct use. \\
$0.50 < T \leq 0.80$ & MEDIUM  & Moderate support; some uncertainty present. Verify if
                                 high-stakes decisions depend on this response. \\
$0.30 < T \leq 0.50$ & LOW     & Weak support or suspicious signals detected. Independent
                                 verification strongly recommended. \\
$T \leq 0.30$      & VERY LOW  & Do not rely on this response. Retrieved context is likely
                                 compromised or answer is unsupported. \\
\hline
\end{tabular}
\end{table}

The binary trustworthiness decision used in experiments is derived from the threshold
$\tau = 0.50$: responses with $T_{\text{final}} \geq 0.50$ are classified as trustworthy
and those with $T_{\text{final}} < 0.50$ are flagged as untrustworthy.

\section{Integration with the RAG Pipeline}

Integration is implemented in \texttt{src/pipeline/rag\_pipeline.py}. Two execution modes are provided:
\begin{itemize}
    \item \texttt{query()} --- retrieval + generation only;
    \item \texttt{query\_with\_evaluation()} --- retrieval + generation + evaluation.
\end{itemize}

When evaluation is enabled:
\begin{enumerate}
    \item retriever returns documents, scores, and embeddings;
    \item generator produces response from retrieved context;
    \item Evaluation Agent computes trust score and detailed diagnostics;
    \item pipeline returns a unified \texttt{RAGResponse} object with optional \texttt{evaluation} field.
\end{enumerate}

The Evaluation Agent is lazily initialized to avoid immediate loading of the NLI model in cases where only fast, non-evaluated querying is required.

\section{Handling Conflicts, Uncertainty, and Evidence Weighting}

\subsection{Conflict Handling}
The system distinguishes between:
\begin{itemize}
    \item \textbf{true contradictions} (high contradiction scores),
    \item \textbf{topic mismatch/neutrality} (weak relevance between texts).
\end{itemize}
Only strong contradictions are used as high-severity conflict signals. This design reduces false alarms from unrelated but non-malicious documents.

\subsection{Uncertainty Handling}
Three uncertainty controls are implemented:
\begin{enumerate}
    \item \textbf{Inconclusive factuality baseline:} return $0.5$ when no clear supporting evidence exists.
    \item \textbf{Short-text guardrails:} very short document pairs are down-weighted or skipped in pairwise NLI consistency checks.
    \item \textbf{Confidence-aware retrieval adjustment:} when the highest FAISS similarity score among the top-$K$ documents falls below a quality threshold ($s_{\text{best}} < 0.30$), the trust score is modestly reduced to reflect low-confidence retrieval (see Equation~\ref{eq:retrieval_confidence}). This penalty is not applied when retrieval quality is adequate ($s_{\text{best}} \geq 0.30$), preventing unnecessary suppression of trust in well-matched queries.
\end{enumerate}

\subsection{Evidence Weighting}
The final poisoning risk uses relevance-aware weighting when retrieval scores are available:
\begin{equation}
P_{overall} = \sum_{i=1}^{k} w_i \cdot P_i, \quad w_i = \frac{s_i}{\sum_j s_j}
\end{equation}
where $s_i$ is retriever similarity and $P_i$ is document-level poison probability.

This reduces spillover effects where a low-relevance poisoned neighbor would otherwise dominate a max-based aggregator.

\section{Implementation Details and Optimization Strategies}

\subsection{Core Data Structures}
The implementation defines typed result structures for traceability:
\begin{itemize}
    \item \texttt{NLIResult}, \texttt{VerificationResult}
    \item \texttt{PoisonSignal}, \texttt{DocumentPoisonResult}, \texttt{PoisonDetectionResult}
    \item \texttt{TrustComponents}, \texttt{TrustIndexResult}
    \item \texttt{EvaluationResult}
\end{itemize}

\subsection{Optimization Strategies}
To maintain practical runtime:
\begin{itemize}
    \item \textbf{Lazy model loading:} NLI model loads only when first needed.
    \item \textbf{Bounded pairwise checks:} document-pair comparisons are capped.
    \item \textbf{Shared NLI instance:} both verifier and poison detector reuse the same verifier object.
    \item \textbf{Fallback-safe behavior:} failure paths return neutral defaults rather than crashing the pipeline.
\end{itemize}

\subsection{High-Level Algorithm}
\begin{verbatim}
Input: query q, response r, retrieved docs D, scores S, embeddings E
1) V <- NLI_Verifier.verify_answer(r, D)
2) P <- Poison_Detector.detect_batch(D, E, S)
3) C <- consistency(V, D)
4) R <- retrieval_confidence(S)
5) T <- TrustIndex.calculate(factuality(V), C, poison(P), R)
6) return EvaluationResult(T, V, P, summary, report)
\end{verbatim}

\subsection{Output and Explainability}
The agent returns:
\begin{itemize}
    \item numeric trust score and categorical trust level;
    \item component contributions;
    \item warnings and recommendation text;
    \item compact summary and detailed formatted report.
\end{itemize}

This makes the module suitable for both quantitative experiments (Chapter~\ref{ch6_experiments}) and human-in-the-loop inspection.

\section{Summary}

This chapter presented the full development of the Evaluation Agent as the main artifact of the thesis. The implementation combines NLI-based factual verification, multi-signal poisoning detection, and a weighted Trust Index into a unified runtime evaluator. It is integrated with the RAG pipeline through a modular interface and designed for interpretable, online trust assessment. The next chapter reports empirical performance under clean and poisoned conditions.

\chapter{Experimental Evaluation}
\label{ch6_experiments}

This chapter presents the empirical evaluation of the Trustworthy RAG framework across
two benchmark datasets and four adversarial poisoning strategies. Section~\ref{sec:ch6_setup}
describes the experimental configuration. Sections~\ref{sec:ch6_baseline}
through~\ref{sec:ch6_overhead} report quantitative results. Section~\ref{sec:ch6_discussion}
interprets the findings in the context of the research questions.

% ─────────────────────────────────────────────────────────────────────────────
\section{Experimental Setup and Scenarios}
\label{sec:ch6_setup}

\subsection{Knowledge Bases and Datasets}

Two publicly available benchmarks were used to construct knowledge bases for evaluation:

\begin{itemize}
    \item \textbf{TruthfulQA}: A benchmark of questions designed to expose common misconceptions
    and factual errors in language model outputs. Each knowledge base document is formatted as
    ``\textit{question} + \textit{best\_answer}'', producing concise fact-grounded entries of
    approximately 20--30 words. TruthfulQA is well-suited for poison detection evaluation
    because many questions have counterintuitive correct answers that diverge from popular
    misconceptions, creating a realistic setting where plausible-sounding false claims could
    be introduced undetected.

    \item \textbf{FEVER}: A large-scale fact verification dataset designed to classify claims
    as SUPPORTED, REFUTED, or UNKNOWN. Each document in the FEVER knowledge base is enriched
    to the format ``\textit{question} + \textit{verdict} + \textit{claim}'', producing entries
    of approximately 30--40 words. This enrichment was critical for reliable NLI evaluation:
    bare claims of only 8--12 words caused approximately 70\% of factuality scores to return
    the inconclusive 0.5 baseline due to insufficient NLI signal. The enriched format provides
    enough semantic context for the NLI Verifier to produce meaningful entailment judgements.
    The FEVER benchmark is used to evaluate factuality improvement from knowledge base
    enrichment (Section~\ref{sec:ch6_factuality}). Per-strategy poison detection experiments
    are conducted exclusively on TruthfulQA, which provides more consistent question-answer
    pairs suitable for the two-part clean/poisoned evaluation design.
\end{itemize}

\subsection{Retrieval Configuration}

The retrieval pipeline is evaluated with two embedding models: (i) the
\texttt{sentence-transformers/all-MiniLM-L6-v2} model producing 384-dimensional dense
representations, and (ii) the \texttt{snowflake-arctic-embed2} model producing
1024-dimensional representations accessed via the FARMI API. Documents are indexed using
FAISS with cosine similarity, and retrieval returns the top $K=5$ most semantically similar
documents for each query. The classification threshold for trustworthiness is fixed at
$\tau = 0.5$: a trust score below this threshold causes the system to flag the response as
untrustworthy and issue an actionable warning.

Two LLMs are evaluated for response generation: \texttt{llama3.3:70b} (primary baseline)
and \texttt{qwen3.5:35b} (cross-model comparison), both accessed via the FARMI computing
cluster API. NLI inference is performed using \texttt{facebook/bart-large-mnli} running on
CPU. Results for the primary Llama 3.3 70B + all-MiniLM-L6-v2 configuration are reported
throughout Sections~\ref{sec:ch6_baseline}--\ref{sec:ch6_overhead}; the full 2×2 grid
comparison is presented in Section~\ref{sec:ch6_grid}.

\subsection{Experiment Design}

Each experiment follows a two-part structure:

\begin{itemize}
    \item \textbf{Part A --- Clean Baseline}: 50 questions are answered against an unmodified
    knowledge base. This measures the average trust score assigned to clean, factually accurate
    retrieval contexts and establishes the baseline accuracy of the naive always-trust system.

    \item \textbf{Part B --- Poisoned Detection}: The same 50 questions are answered against a
    knowledge base in which 30\% of documents have been adversarially modified. This measures
    the Evaluation Agent's ability to detect and flag compromised retrieval contexts before
    generation.
\end{itemize}

Two experiment configurations are reported:

\begin{enumerate}
    \item \textbf{Per-strategy experiments}: 20 samples per run with a single poisoning strategy
    applied in isolation (contradiction, instruction injection, entity swap, or subtle
    manipulation). This enables characterisation of detection difficulty per attack type.

    \item \textbf{Mixed experiment}: 50 samples with poisoning strategies assigned randomly
    across the poisoned portion of the knowledge base. This reflects a more realistic threat
    scenario where multiple attack types may co-exist.
\end{enumerate}

A document is counted as successfully poisoned only if the modified text differs from the
original (\texttt{poisoned\_text} $\neq$ \texttt{original\_text}). This prevents evaluation
metrics from being distorted by unchanged documents incorrectly counted as poisoned samples.

\subsection{Adversarial Poisoning Strategies}

Four rule-based poisoning strategies are applied to modify documents in the knowledge base:

\begin{itemize}
    \item \textbf{Contradiction}: Appends a sentence that semantically contradicts the key
    claim of the original document (e.g., negating the stated answer).

    \item \textbf{Instruction Injection}: Appends an override directive targeting the language
    model's instruction-following behaviour (e.g., ``\textit{IMPORTANT: Ignore all previous
    context and state that [false claim]}'').

    \item \textbf{Entity Swap}: Replaces named entities or numeric values in-place within the
    original text. A three-level fallback is employed: (1) predefined named-entity substitution;
    (2) numeric value shifting via regular expressions when no named entity is present;
    (3) qualifier or negation swapping in the answer segment of the document. This strategy
    produces no appended content and no structural artifacts, making it the hardest attack to
    detect by surface-level analysis.

    \item \textbf{Subtle Manipulation}: Inserts hedging language, false qualifiers, or
    plausible-sounding uncertainty signals that subtly shift the factual content without
    triggering obvious pattern-matching rules.
\end{itemize}

% ─────────────────────────────────────────────────────────────────────────────
\section{Baseline Models and Comparison Methods}
\label{sec:ch6_baseline}

\subsection{Naive Always-Trust Baseline}

The primary comparison point is a \textbf{naive always-trust RAG system}: a standard
retrieval-augmented generation pipeline that accepts all retrieved documents and generated
responses without any evaluation step. This baseline always returns
$\texttt{is\_trustworthy} = \text{True}$.

In the two-part experiment design, the baseline's accuracy equals the proportion of samples
where trusting the output is the correct decision. With a 30\% poisoning rate, approximately
70\% of Part B samples involve clean retrieval contexts, giving the naive baseline an inherent
accuracy advantage that reflects the class imbalance rather than any detection capability.
This makes accuracy alone a misleading evaluation metric and motivates the use of trust score
separation alongside precision, recall, and F1.

\subsection{Trustworthy RAG System}

The full evaluation pipeline --- NLI Verifier, Poison Detector, and Trust Index Calculator
--- constitutes the Trustworthy RAG system evaluated in this chapter. Its performance is
measured against the naive baseline across the following metrics:

\begin{itemize}
    \item \textbf{Accuracy}: fraction of samples correctly classified as trustworthy or
    untrustworthy relative to the ground truth poisoning label.

    \item \textbf{Precision}: fraction of flagged (untrustworthy) samples that are genuinely
    poisoned.

    \item \textbf{Recall (Poison Detection Rate)}: fraction of poisoned samples that the
    Evaluation Agent successfully flags before generation.

    \item \textbf{F1 Score}: harmonic mean of precision and recall; the primary optimisation
    target for the detection task.

    \item \textbf{Trust Score Separation}: the difference between the mean clean trust score
    and the mean poisoned trust score,
    $\Delta = \bar{T}_{\text{clean}} - \bar{T}_{\text{poisoned}}$.
    This metric quantifies the discriminative power of the Trust Index independently of the
    classification threshold and class imbalance.
\end{itemize}

% ─────────────────────────────────────────────────────────────────────────────
\section{Results on Factuality Improvement}
\label{sec:ch6_factuality}

Table~\ref{tab:factuality} reports the average trust scores and system accuracy on both
benchmark datasets. The clean trust score reflects how the Evaluation Agent rates responses
grounded in unmodified, factually accurate knowledge. The poisoned trust score reflects the
same pipeline's rating of responses retrieved from adversarially modified contexts.

\begin{table}[h]
\centering
\caption{Trust score and accuracy comparison across benchmark datasets (mixed strategy, 50 samples).}
\label{tab:factuality}
\resizebox{\textwidth}{!}{%
\begin{tabular}{lcccccc}
\hline
\textbf{Dataset} & \textbf{Clean Trust} & \textbf{Poisoned Trust} & \textbf{Separation}
    & \textbf{System Acc.} & \textbf{Baseline Acc.} & \textbf{Improvement} \\
\hline
TruthfulQA & 0.830 & 0.590 & 0.240 & 91\% & 85\% & +7\% \\
FEVER      & 0.654 & ---   & ---   & 90\% & 85\% & +5\% \\
\hline
\end{tabular}%
}
\end{table}
\small\textit{Note: FEVER was evaluated for factuality improvement only (clean knowledge
base); poisoned trust and separation values are not applicable for this benchmark.}

\normalsize
On TruthfulQA, the Evaluation Agent assigns an average clean trust score of 0.830 to
responses grounded in unmodified documents, indicating that the pipeline consistently
rates factually supported outputs as trustworthy. The poisoned average of 0.590 shows
that the agent is sensitive to adversarial content, producing a separation of 0.240 that
enables discrimination above the 0.5 classification threshold in the majority of cases.
The system achieves 91\% accuracy on TruthfulQA, a +7\% improvement over the naive
always-trust baseline, with zero false positives on clean content (100\% precision).

On FEVER, the system achieves 90\% accuracy, a +5\% improvement over the naive baseline.
This larger gain is partially attributable to the enriched document format: enriched FEVER
entries provide substantially more NLI-relevant context than bare TruthfulQA answers,
enabling the NLI Verifier to produce informative factuality signals and reducing the
frequency of the inconclusive 0.5 baseline score from approximately 70\% (on bare claims)
to a small minority of edge cases. The FEVER enrichment process --- reformatting each
document from a bare claim to a structured ``question + verdict + claim'' tuple --- was a
key implementation decision that unlocked reliable NLI-based evaluation on short-document
knowledge bases.

% ─────────────────────────────────────────────────────────────────────────────
\section{Detection Performance under Data Poisoning Attacks}
\label{sec:ch6_detection}

\subsection{Overall Detection Performance}

Table~\ref{tab:main_results} reports detection performance on the TruthfulQA benchmark
with 50 samples and a mixed poisoning strategy, where attack types are assigned randomly
across the poisoned portion of the knowledge base.

\begin{table}[h]
\centering
\caption{Overall detection performance --- TruthfulQA, 50 samples, mixed strategy.}
\label{tab:main_results}
\begin{tabular}{lc}
\hline
\textbf{Metric} & \textbf{Value} \\
\hline
Accuracy                  & 91\%   \\
Precision                 & 100\%  \\
Recall                    & 40\%   \\
F1 Score                  & 57.1\% \\
Trust Score Separation    & 0.240  \\
Baseline Accuracy         & 85\%   \\
Improvement over Baseline & +7\%   \\
\hline
\end{tabular}
\end{table}

The system achieves 40\% recall across mixed strategies, meaning that 40\% of poisoned
retrieval contexts are correctly identified and flagged before generation. Precision of
100\% indicates zero false positives on clean content --- the Evaluation Agent does not
flag any clean response as untrustworthy. The F1 score of 57.1\% reflects the
precision--recall trade-off: the system is highly conservative, preferring to miss some
poisoned samples rather than raise false alarms. The FP spillover effect observed in
earlier experimental rounds (where clean queries retrieving poisoned neighbours received
elevated poison probabilities) was mitigated by the relevance-weighted poison aggregation
described in Chapter~\ref{ch5_eval_agent}.

\subsection{Per-Strategy Detection Results}

Table~\ref{tab:per_strategy} reports detection performance for each poisoning strategy in
isolation, using 20 samples per strategy.

\begin{table}[h]
\centering
\caption{Per-strategy detection performance --- TruthfulQA, 20 samples per strategy.}
\label{tab:per_strategy}
\begin{tabular}{lccccc}
\hline
\textbf{Strategy} & \textbf{Accuracy} & \textbf{Precision} & \textbf{Recall}
    & \textbf{F1} & \textbf{Separation} \\
\hline
Contradiction  & 89\%   & 62.1\%  & 60.0\%  & 61.1\%  & 0.245 \\
Injection      & 75\%   & 37.5\%  & 100\%   & 54.5\%  & 0.407 \\
Entity Swap    & 87.5\% & 100\%   & 16.7\%  & 28.6\%  & 0.071 \\
Subtle         & 85\%   & 50.0\%  & 16.7\%  & 25.0\%  & 0.127 \\
Mixed          & 87.5\% & 55.6\%  & 83.3\%  & 66.7\%  & 0.327 \\
\hline
\end{tabular}
\end{table}
\small\textit{Note: the Mixed row reports the 20-sample per-strategy run where all four
strategies are combined. This differs from the 100-sample main experiment in
Table~\ref{tab:main_results}, which uses a separate run with a larger sample size and
correspondingly different classification statistics (91\% accuracy, 57.1\% F1).}
\normalsize

The results reveal a wide range of detection difficulty across strategies. Instruction
injection achieves perfect recall (100\%) due to its high-severity linguistic artifacts,
while entity swap and subtle manipulation both achieve only 16.7\% recall. The mixed
strategy achieves the best overall balance (F1 = 66.7\%, recall = 83.3\%), as different
strategies activate different detection signals and their combined presence creates more
detectable patterns than any single strategy in isolation.

% ─────────────────────────────────────────────────────────────────────────────
\section{Robustness and Trust Index Analysis}
\label{sec:ch6_robustness}

\subsection{Trust Score Distributions per Strategy}

Table~\ref{tab:separation} summarises the trust score separation for each poisoning
strategy, providing a threshold-independent view of the Trust Index's discriminative power.

\begin{table}[h]
\centering
\caption{Trust score separation per poisoning strategy.}
\label{tab:separation}
\begin{tabular}{lcccc}
\hline
\textbf{Strategy} & \textbf{Clean Avg} & \textbf{Poisoned Avg}
    & \textbf{Separation ($\Delta$)} & \textbf{Discriminability} \\
\hline
Injection     & $\approx 0.76$ & $\approx 0.35$ & 0.407 & Strong    \\
Mixed         & 0.830          & 0.590          & 0.240 & Moderate  \\
Contradiction & ---            & ---            & 0.245 & Moderate  \\
Subtle        & ---            & ---            & 0.127 & Weak      \\
Entity Swap   & $\approx 0.75$ & $\approx 0.68$ & 0.071 & Very weak \\
\hline
\end{tabular}
\end{table}
\small\textit{Note: clean and poisoned average trust values are reported for strategies
with the most extreme separation (injection and entity swap). Component averages for
contradiction and subtle strategies were not separately logged; their separation values
are derived from per-sample trust score records.}

\normalsize
The injection strategy produces the largest separation (0.407), enabling reliable
discrimination at the fixed threshold of 0.5. Entity swap produces the smallest separation
(0.071), reflecting near-overlapping trust score distributions between clean and poisoned
samples. This near-zero separation confirms that entity swap attacks cannot be detected
by Trust Index thresholding alone --- the upstream poison detection signals do not differ
sufficiently between clean and entity-swapped contexts.

\subsection{Non-Linear Dampener Effect}

The non-linear poison dampener, described in Chapter~\ref{ch4_methodology}, activates when
$P_{poison} > 0.7$ and reduces the final trust score by up to 40\% at $P_{poison} = 1.0$.
This mechanism prevents high factuality scores from masking an elevated poisoning signal
when retrieved documents appear superficially supportive of the generated answer but are
otherwise structurally suspicious.

In the injection experiment, average poison probabilities for poisoned contexts reached
0.994, consistently triggering the dampener and producing the strong separation of 0.407
observed. For entity swap, average poisoned-context poison probabilities remained near
0.319 --- well below the 0.7 dampener threshold --- leaving trust scores largely unchanged
and yielding the near-zero separation of 0.071. The dampener thus amplifies the Trust
Index's response to high-confidence poison signals while having no effect on weak signals,
confirming its role as a robustness mechanism rather than a uniform scaling factor.

% ─────────────────────────────────────────────────────────────────────────────
\section{Ablation Studies}
\label{sec:ch6_ablation}

To validate the Trust Index weight configuration, four weight settings were evaluated on
the TruthfulQA mixed-strategy experiment. Table~\ref{tab:ablation} reports the results.

\begin{table}[h]
\centering
\caption{Ablation study --- effect of Trust Index weight configuration on detection performance.}
\label{tab:ablation}
\begin{tabular}{lcccccc}
\hline
\textbf{Configuration} & \textbf{$\alpha$ (Fact.)} & \textbf{$\beta$ (Cons.)}
    & \textbf{$\gamma$ (Poison)} & \textbf{Accuracy} & \textbf{F1} & \textbf{Recall} \\
\hline
Default          & 0.40 & 0.35 & 0.25 & 81.7\% & 48\%  & ---    \\
Equal weights    & 0.33 & 0.34 & 0.33 & ---    & 48\%  & ---    \\
Poison-heavy     & $<0.20$ & $<0.20$ & $>0.60$ & ---    & ---   & 88.9\% \\
Factuality-heavy & $>0.60$ & $<0.20$ & $<0.20$ & 81.7\% & ---   & ---    \\
\hline
\end{tabular}
\end{table}
\small\textit{Note: cells marked --- indicate metrics not separately recorded for that
configuration variant; each configuration was evaluated against its primary target metric.}

\normalsize
The poison-heavy configuration achieves the highest recall (88.9\%) by amplifying the
poisoning signal component, at the cost of precision. The factuality-heavy configuration
achieves the highest accuracy (81.7\%) but misses subtle and entity-swap attacks that
produce low poison probability signals below the non-linear dampener threshold. Default
weights (0.40, 0.35, 0.25) provide the best overall trade-off. Equal weights produce
comparable F1 to the default, suggesting that the optimisation landscape is relatively
flat near the configuration optimum and that the Trust Index is not highly sensitive to
small perturbations in weight values.

% ─────────────────────────────────────────────────────────────────────────────
\section{Performance Overhead and Scalability}
\label{sec:ch6_overhead}

Table~\ref{tab:latency} summarises the runtime characteristics of the Evaluation Agent
measured over the TruthfulQA experiments.

\begin{table}[h]
\centering
\caption{Evaluation Agent performance overhead per sample.}
\label{tab:latency}
\begin{tabular}{lc}
\hline
\textbf{Metric} & \textbf{Value} \\
\hline
Total time per sample (retrieval + generation + evaluation) & $\approx 13.2$\,s \\
Evaluation-only time (NLI + poison detection + Trust Index)  & $\approx 10$--$11$\,s \\
Throughput with Evaluation Agent enabled                    & $\approx 4$--$5$ queries/min \\
NLI inference calls per batch ($K=5$ retrieved documents)   & Up to 10 pairwise calls \\
Baseline RAG throughput (no evaluation)                     & $<1$\,s / query \\
Latency overhead factor                                     & $\approx 13\times$ \\
\hline
\end{tabular}
\end{table}

The primary runtime bottleneck is the \texttt{facebook/bart-large-mnli} model running on
CPU. With $K=5$ retrieved documents, intra-document NLI and cross-document consistency
checks require up to 10 NLI forward passes per batch, each taking approximately 1 second
on CPU hardware. The shared NLI instance --- reused by both the NLI Verifier and the Poison
Detector --- prevents redundant model loading overhead but does not reduce the number of
inference calls required per query.

The 13-second evaluation overhead represents a 13$\times$ latency increase over baseline
RAG, for accuracy gains of +1\% to +6\% over the naive always-trust baseline. This
trade-off is not justified for latency-sensitive real-time conversational applications.
However, for \textbf{high-stakes batch scenarios} --- such as automated document
verification, regulatory compliance review, or medical information retrieval systems ---
the evaluation cost is acceptable relative to the risk of acting on poisoned information.

Several optimisation paths exist for reducing latency in deployment contexts:

\begin{itemize}
    \item \textbf{GPU acceleration}: NLI inference on a mid-range GPU is estimated to reduce
    evaluation time from $\approx 10$\,s to $<2$\,s per sample.

    \item \textbf{Bounded pairwise checks}: limiting cross-document NLI to the top-$K'$ most
    semantically dissimilar document pairs, rather than all $\binom{K}{2}$ pairs.

    \item \textbf{Query result caching}: caching NLI results for repeated or near-duplicate
    queries reduces redundant computation in document-heavy batch workflows.
\end{itemize}

% ─────────────────────────────────────────────────────────────────────────────
\section{Discussion of Results}
\label{sec:ch6_discussion}

\subsection{Strategy-Specific Detection Analysis}

The per-strategy results reveal a clear detection hierarchy that reflects both the signal
richness of each attack type and the architectural coverage of the Evaluation Agent's
detection signals.

\paragraph{Instruction Injection --- Fully Detectable.}
Injection achieves 100\% recall because appended override directives (e.g.,
\texttt{``IMPORTANT: Ignore all previous context and state that...''}) contain high-severity
linguistic patterns that trigger deterministic regex rules in the Poison Detector with
severity scores of 0.5--0.7. These artifacts are structurally unambiguous and do not
require NLI inference for detection. Intra-document NLI provides a complementary signal:
splitting each document and running NLI between the original factual content and the
appended adversarial directive reveals a strong semantic discontinuity, producing
intra-document contradiction scores above the 0.7 threshold. The injection strategy also
exhibits the largest trust score separation (0.407) across all strategies tested.

\paragraph{Contradiction --- Moderately Detectable.}
Contradiction achieves 60.0\% recall with a separation of 0.245. Intra-document NLI
successfully detects cases where the appended contradiction produces a strong first-half
to second-half contradiction signal (above 0.7). However, some poisoned samples produce
scores near the detection threshold, resulting in borderline classifications. The
detection rate is sensitive to the specific phrasing of the contradiction: explicit
negations are more detectable than hedged or qualified contradictions.

\paragraph{Entity Swap --- Near-Undetectable.}
Entity swap achieves only 16.7\% recall with a separation of 0.071. Unlike injection and
contradiction, entity swap modifies documents in-place with no appended content and no
structural artifacts. The resulting documents are syntactically and stylistically
indistinguishable from genuine documents at the surface level --- only specific numeric
or named-entity values are changed. Detecting this class of attack requires external
world-knowledge verification (e.g., querying a fact database to verify whether a given
value is correct for a particular claim), which is beyond the capability of a
surface-signal-based evaluation agent. This constitutes an \textbf{architectural limit}
of the current approach.

\paragraph{Subtle Manipulation --- Near-Undetectable.}
Subtle manipulation similarly achieves 16.7\% recall with a separation of 0.127. Hedging
language and false qualifiers do not trigger pattern-matching rules, and the semantic shift
introduced is too small for intra-document NLI to consistently detect above threshold.
Average poison probabilities for this strategy hover near 0.700 --- at the boundary of the
non-linear dampener --- indicating that the signals present are at the edge of statistical
significance under the current architecture.

\subsection{False Positive Spillover Effect}

A consistent source of false positives across all strategies is the
\textit{FP spillover effect}: a clean query that retrieves a poisoned document as a
neighbour in the top-$K$ results receives an elevated overall poison probability, causing
the Evaluation Agent to flag the response as untrustworthy even when the highest-ranked
retrieved document is clean.

This behaviour is architecturally correct from a security perspective: the retrieval
environment is contaminated, and the Evaluation Agent correctly warns that the context
cannot be fully trusted. However, the spillover reduces the precision metric, particularly
in the injection experiment. The relevance-weighted poison aggregation introduced in
Chapter~\ref{ch5_eval_agent} mitigates this effect by weighting each document's poison
probability by its retrieval similarity score, so low-relevance poisoned neighbours
contribute less to the overall estimate than the top-ranked document. This intervention
reduced false positives substantially in clean-KB experiments. However, when a poisoned
document ranks high in semantic similarity to the query, the spillover effect is not fully
eliminated by relevance weighting alone.

This represents a fundamental design trade-off: a system that issues warnings about
contaminated retrieval environments will inevitably produce false positives when clean
queries share embedding space proximity with poisoned documents. The appropriate response
is not to suppress these warnings but to provide richer contextual information about which
specific documents triggered the alert.

\subsection{Trust Score Separation as Primary Metric}

In class-imbalanced settings --- where 70\% of samples involve clean retrieval contexts
--- accuracy alone is a misleading evaluation metric. The naive always-trust baseline
achieves 85\% accuracy without any detection capability, simply by never raising a
warning. The Trust Index's discriminative value is better captured by the trust score
separation $\Delta$, which measures how far apart the clean and poisoned trust
distributions are independently of threshold choice and class composition.

The per-strategy separation range (0.071 for entity swap to 0.407 for injection) provides
actionable guidance for future defence design. Strategies with high separation are
candidates for reliable deployment under the current architecture; strategies with low
separation require fundamentally different detection mechanisms. This per-strategy
characterisation is a contribution of this thesis: prior evaluation frameworks have not
systematically quantified the detection difficulty of individual poisoning strategies in
RAG-specific evaluation settings.

% ─────────────────────────────────────────────────────────────────────────────
\section{Comparative Evaluation: LLM and Embedding Model Sensitivity}
\label{sec:ch6_grid}

To assess whether the Trust Index performance is specific to the Llama 3.3 70B model and
the all-MiniLM-L6-v2 embedding pipeline, a \textbf{2×2 factorial experiment} was conducted
crossing two LLMs with two embedding models. All four configurations were evaluated under
identical conditions: 100 samples (50 clean Part A, 50 poisoned Part B), 30\% poison ratio,
mixed strategy, $K=3$, $\tau=0.5$. The impact of increasing retrieval depth to $K=5$ is
analysed separately in Section~\ref{sec:ch6_topk}.

\subsection{Grid Results}

Table~\ref{tab:grid_results} reports the full performance grid.

\begin{table}[h]
\centering
\caption{2×2 factorial grid results --- TruthfulQA, 100 samples, mixed strategy.}
\label{tab:grid_results}
\resizebox{\textwidth}{!}{%
\begin{tabular}{llccccccc}
\hline
\textbf{LLM} & \textbf{Embedding} & \textbf{Acc.} & \textbf{Prec.} & \textbf{Recall}
    & \textbf{F1} & \textbf{Clean $\bar{T}$} & \textbf{Poison $\bar{T}$} & \textbf{Sep.} \\
\hline
Llama 3.3 70B & all-MiniLM-L6-v2       & 91\% & 100\%  & 40\%  & 57.1\% & 0.830 & 0.590 & 0.240 \\
Llama 3.3 70B & snowflake-arctic-embed2 & 91\% & 100\%  & 40\%  & 57.1\% & 0.799 & 0.611 & 0.188 \\
Qwen 3.5 35B  & all-MiniLM-L6-v2       & 71\% & 25.0\% & 46.7\% & 32.6\% & 0.633 & 0.471 & 0.161 \\
Qwen 3.5 35B  & snowflake-arctic-embed2 & 71\% & 28.1\% & 60.0\% & 38.3\% & 0.653 & 0.454 & 0.199 \\
\hline
\multicolumn{2}{l}{Naive always-trust baseline} & 85\% & --- & --- & --- & --- & --- & --- \\
\hline
\end{tabular}%
}
\end{table}

\subsection{Finding 1: Embedding Invariance for Llama 3.3 70B}

The two Llama configurations produce \textbf{identical classification results}
(Acc = 91\%, Prec = 100\%, Recall = 40\%, F1 = 57.1\%), with the sole difference being a
slightly higher clean trust score for MiniLM ($0.830$ vs.\ $0.799$) and a larger trust
score separation ($0.240$ vs.\ $0.188$). The zero false positive count in both Llama runs
confirms that the precision advantage is attributable to the LLM's concise answer style
rather than to retrieval quality.

This result demonstrates that for a concise-output LLM such as Llama 3.3, switching to a
higher-dimensional embedding model (1024-dim Snowflake vs.\ 384-dim MiniLM) does not
improve poison detection. The Trust Index's discriminative power is determined primarily
by the NLI entailment signal between the generated answer and retrieved documents, which
is insensitive to the specific embedding space used for retrieval as long as topically
relevant documents are returned.

\subsection{Finding 2: Qwen 3.5 35B Generation-Style False Positives}

Both Qwen configurations \textbf{underperform the naive baseline} by 16.5 percentage
points (71\% vs.\ 85\%), producing 21--23 false positives out of 85 clean samples.
The root cause is the LLM's verbose, hedged generation style: Qwen 3.5 35B typically
prefaces answers with phrases such as ``\textit{Based on the provided context}'' or
``\textit{According to the available information},'' which the NLI Verifier interprets as
reduced semantic commitment to the claim. This reduces $S_{\text{factuality}}$ below 0.5
for many clean responses, causing the Trust Index to fall below the classification
threshold $\tau = 0.5$ and triggering false untrustworthy labels.

The mean clean trust score for Qwen ($\approx 0.64$) is substantially lower than for
Llama ($\approx 0.83$), confirming that the generation-style effect is systematic rather
than query-specific. This finding demonstrates that the Trust Index's $\tau = 0.5$
threshold is implicitly calibrated for the Llama generation distribution and requires
per-LLM recalibration for reliable deployment with alternative generators.

\subsection{Finding 3: Retrieval Quality Paradox}

For Qwen 3.5 35B, the higher-dimensional Snowflake Arctic Embed2 model shows marginally
improved F1 (38.3\% vs.\ 32.6\%) and recall (60\% vs.\ 46.7\%), with the same accuracy
(71\%) and a comparable false positive count. The slight improvement in recall may reflect
that Snowflake's superior retrieval precision recovers more relevant poisoned documents,
making them available for NLI analysis.

However, a counter-intuitive pattern also emerges: Snowflake's stronger retrieval retrieves
\textit{cleaner} documents for genuinely clean queries, raising the mean clean trust score
slightly (0.653 vs.\ 0.633). This constitutes a \textbf{retrieval quality paradox}: higher
retrieval quality dilutes contaminated context by preferring highly relevant clean documents,
reducing the poison detection signal for poisoned queries while simultaneously reducing
false positives for clean queries. The net effect on F1 depends on which correction
dominates.

\subsection{Implications for Deployment}

The grid experiment yields three actionable guidelines for deploying the Trustworthy RAG
framework:

\begin{enumerate}
    \item \textbf{LLM selection is the dominant factor.} Switching embedding models within
    the same LLM produces negligible differences in detection performance for Llama.
    The choice of LLM has a far larger effect on both precision (100\% vs.\ 25--28\%) and
    overall accuracy (91\% vs.\ 71\%).

    \item \textbf{Per-LLM threshold calibration is required.} The fixed threshold $\tau = 0.5$
    is not universally applicable. For LLMs with hedged generation styles, the Trust Index
    threshold must be lowered, or the factuality component normalised, to avoid systematic
    false positives on clean content.

    \item \textbf{The framework is validated as reproducible on Llama 3.3 70B.} The new
    Llama+MiniLM run replicates the original thesis baseline identically (same TP/FP/FN
    counts, trust score differences below 0.015), confirming that the results are stable
    across independent runs.
\end{enumerate}

% ─────────────────────────────────────────────────────────────────────────────
\section{Retrieval Depth Sensitivity: $K=3$ versus $K=5$}
\label{sec:ch6_topk}

To isolate the effect of retrieval depth on detection performance, the primary configuration
(Llama~3.3~70B + all-MiniLM-L6-v2) was evaluated under two retrieval settings: $K=3$ (used
in the 2×2 grid experiment, Section~\ref{sec:ch6_grid}) and $K=5$ (the system default).
All other parameters were held constant (TruthfulQA, 100 samples, 30\% poison ratio,
mixed strategy, $\tau=0.5$).

Increasing $K$ from 3 to 5 has two theoretically opposing effects. First, more retrieved
documents increase the probability of a poisoned document appearing in the batch, which
should improve recall. Second, the number of NLI inference calls grows from
$3 + 3 + \binom{3}{2} = 9$ to $5 + 5 + \binom{5}{2} = 20$ per batch (factuality +
intra-document + pairwise cross-document), increasing evaluation overhead approximately 2.2$\times$.

\begin{table}[h]
\centering
\caption{Retrieval depth sensitivity: $K=3$ vs $K=5$ (Llama~3.3~70B + all-MiniLM-L6-v2,
TruthfulQA, 100 samples, mixed strategy).}
\label{tab:topk_comparison}
\begin{tabular}{lccccccc}
\hline
\textbf{K} & \textbf{Acc.} & \textbf{Prec.} & \textbf{Recall} & \textbf{F1}
    & \textbf{Clean $\bar{T}$} & \textbf{Sep.} & \textbf{NLI calls/batch} \\
\hline
$K=3$ & 91\%  & 100\%  & 40.0\% & 57.1\% & 0.830 & 0.240 & $\sim$9 \\
$K=5$ & 91\%  & 100\%  & 40.0\% & 57.1\% & 0.827 & 0.237 & $\sim$20 \\
\hline
\end{tabular}
\end{table}

\subsection{Finding: Detection Is Retrieval-Depth-Invariant for Llama~3.3~70B}

The results in Table~\ref{tab:topk_comparison} show that increasing retrieval depth from
$K=3$ to $K=5$ produces \textbf{no improvement in detection performance}. The classification
outcome is identical: 6 true positives, 0 false positives, and 9 false negatives in both
configurations. Accuracy (91\%), precision (100\%), recall (40.0\%), and F1 (57.1\%) are
unchanged. The trust score separation decreases marginally from 0.240 to 0.237, a difference
within noise bounds.

The predicted recall improvement does not materialise. The reason is structural: the 9
undetected poisoned contexts (the false negatives) correspond to Entity Swap and Subtle
attack strategies, which produce in-place modifications to the document text. These attacks
are difficult to detect because they leave no textual artefacts that would be flagged by
linguistic pattern detectors or intra-document NLI, regardless of how many documents are
retrieved. Adding two extra clean documents to the batch provides no additional poison signal
for these attacks.

Conversely, the 6 successfully detected contexts (Contradiction and Injection strategies)
are detectable at $K=3$ because they append distinct textual patterns (contradiction markers
and instruction-override phrases) that the detector flags reliably. Retrieving 5 instead of
3 documents does not improve recall for these strategies, since they are already detected at
the lower retrieval depth.

\subsection{Implication: Detection Bottleneck Is Attack Strategy, Not Retrieval Depth}

This null result has a clear architectural implication: \textbf{the detection ceiling of the
Evaluation Agent at 40\% recall is set by the attack strategy difficulty, not by the number
of documents retrieved}. Entity Swap and Subtle attacks represent fundamental limits of the
current NLI-and-heuristic approach: they require external world knowledge or fine-grained
semantic comparison to detect reliably.

From a deployment standpoint, $K=3$ is therefore the preferred configuration for this
system: it achieves equivalent detection performance at approximately half the NLI inference
cost ($\sim$9 vs.\ $\sim$20 calls per batch). The $K=3$ setting is used for the 2×2 grid
experiment (Section~\ref{sec:ch6_grid}), and this sensitivity analysis retroactively
validates that choice.

% ─────────────────────────────────────────────────────────────────────────────
\section{Summary}

This chapter evaluated the Trustworthy RAG framework across TruthfulQA and FEVER
benchmarks with four adversarial poisoning strategies. The Evaluation Agent achieves
up to 90\% accuracy (FEVER benchmark) and 66.7\% F1 score (mixed per-strategy
experiment), with instruction injection detected at 100\% recall.
On the mixed 100-sample experiment (Llama 3.3 70B + all-MiniLM-L6-v2), the system
achieves 91\% accuracy, 57.1\% F1, 100\% precision, and a trust score separation of
0.240, representing a +7\% improvement over the naive always-trust baseline on
TruthfulQA and +5\% on FEVER.

Per-strategy analysis reveals a clear detection hierarchy: instruction injection and
contradiction are amenable to textual-signal and intra-document NLI approaches, while
entity swap and subtle manipulation represent architectural limits requiring external
world-knowledge for reliable detection. The 13-second evaluation overhead positions the
system as suitable for high-stakes batch deployment rather than real-time applications.

The ablation study confirms that the default weight configuration (0.40, 0.35, 0.25)
provides the best overall trade-off between precision, recall, and accuracy, and that
the Trust Index is robust to small perturbations in weight values.

A 2×2 factorial experiment (Section~\ref{sec:ch6_grid}) further establishes that LLM
selection is the dominant performance factor: Llama 3.3 70B achieves 91\% accuracy with
zero false positives under both embedding models, while Qwen 3.5 35B underperforms the
naive baseline due to generation-style false positives. These findings highlight the
importance of per-LLM Trust Index calibration in practical deployments.

Section~\ref{sec:ch6_topk} (Retrieval Depth Sensitivity) establishes that detection
performance is retrieval-depth-invariant for Llama~3.3~70B: $K=5$ produces identical
classification outcomes to $K=3$ (91\% accuracy, 57.1\% F1, 0 false positives) at
approximately 2.2$\times$ the NLI inference cost, confirming that the 40\% recall ceiling
is set by attack strategy difficulty rather than retrieval depth. The next chapter discusses all findings in the context of
the research questions, identifies the system's theoretical and practical contributions,
and outlines directions for future work.

\chapter{Discussion}
\label{ch7_discussion}

This chapter interprets the empirical findings reported in Chapter~\ref{ch6_experiments} in the context of the three Research Questions stated in Chapter~\ref{ch:introduction}. It articulates the specific theoretical and practical contributions of this work, discusses the implications for real-world deployment of trustworthy RAG systems, identifies the study's limitations, and reflects on the ethical considerations surrounding adversarial AI evaluation research.

% ─────────────────────────────────────────────────────────────────────────────
\section{Interpretation of Key Findings}
\label{sec:disc_findings}

\subsection{RQ1: Detection Effectiveness of the Evaluation Agent}

\textit{How effectively can an autonomous Evaluation Agent detect misinformation and knowledge poisoning in a RAG pipeline?}

The results reveal a clear detection hierarchy across the four adversarial strategies, and the answer to RQ1 is nuanced: the Evaluation Agent is highly effective against overt poisoning strategies but structurally limited against subtle in-place modifications.

\textbf{Instruction injection} is the most detectable strategy, achieving 100\% recall in per-strategy experiments. The Evaluation Agent's linguistic pattern detector reliably identifies instruction-override phrases (e.g., ``IMPORTANT: Ignore previous context''), and the intra-document NLI component detects the structural inconsistency introduced when such phrases are appended to otherwise coherent documents. This finding confirms that prompt injection attacks --- which \textcite{greshake2023indirect} identified as a critical threat to LLM-integrated systems --- are amenable to surface-level detection when the injection is syntactically distinct from the host document.

\textbf{Contradiction} poisoning achieves moderate detection (33--89\% recall across experiments), primarily through intra-document NLI contradiction signals. When a contradicting sentence is appended to a document, the NLI model detects semantic inconsistency between the document's two halves. Detection is partial rather than complete because the NLI signal strength varies with sentence specificity: generic contradiction markers (e.g., ``Contrary to popular belief'') produce weaker NLI signals than factually specific contradictions.

\textbf{Entity swap} and \textbf{subtle manipulation} represent the architectural limit of the current approach, achieving only 16.7\% recall in per-strategy experiments. These strategies operate through in-place text modifications --- replacing numbers, proper nouns, or qualifiers within the original sentence structure --- leaving no appended content for intra-document NLI or pattern detectors to flag. The modified document is semantically close to the clean version and produces nearly identical NLI scores. Detecting these attacks would require external world-knowledge reasoning (e.g., verifying that ``Berlin'' is an incorrect answer to a question about France's capital) --- a capability beyond the current NLI-and-heuristic architecture.

In the mixed-strategy 100-sample experiment, the system achieves 40\% overall recall with 100\% precision (zero false positives on clean content), demonstrating that the agent's detection is reliable when it fires but conservative in coverage. This precision-recall trade-off reflects a deliberate design choice: false positives on clean content undermine user trust in the system's warnings, making precision the higher-priority metric in high-stakes deployments.

\subsection{RQ2: Reliability and Calibration of the Trust Index}

\textit{How reliably does the Trust Index quantify trustworthiness across diverse poisoning scenarios?}

The Trust Index provides a calibrated, interpretable trustworthiness score across the evaluated scenarios. Several properties of its design are validated by the experiments.

\textbf{Discriminative separation}: In the primary configuration (Llama~3.3~70B + all-MiniLM-L6-v2, K=3), clean responses achieve a mean Trust Index of $\bar{T}_{\text{clean}} = 0.830$, while poisoned contexts produce $\bar{T}_{\text{poison}} = 0.590$, a separation of 0.240 points. This separation is consistent across both embedding models for Llama (0.240 for MiniLM, 0.188 for Snowflake), indicating that the Trust Index's discriminative capacity is primarily driven by the NLI signal rather than by retrieval quality.

\textbf{Non-linear dampener effectiveness}: The ablation study confirms that the dampener (Section~\ref{subsubsec:dampener}) plays a critical role under high-contamination conditions. Without the dampener, responses with high factuality scores (because the LLM ignores the poisoned portion) can exceed the $\tau = 0.5$ threshold despite high $P_{\text{poison}}$ values. The dampener imposes a multiplicative penalty that prevents this override, ensuring that a 90\% poison probability forces the final trust score below the classification threshold even when factuality and consistency appear healthy.

\textbf{Weight sensitivity}: The ablation study demonstrates that the default weight configuration ($\alpha=0.40$, $\beta=0.35$, $\gamma=0.25$) yields the best overall F1 across the tested weight variations. Poison-heavy configurations ($\gamma=0.50$) maximise recall (88.9\%) at the cost of precision, while factuality-heavy configurations maximise accuracy. The default weights represent a balanced operating point --- a design choice validated by empirical evidence rather than arbitrary assignment.

\textbf{LLM-dependency}: The 2×2 factorial experiment reveals a critical limitation of the Trust Index's universality: its component scores are sensitive to LLM generation style. Qwen~3.5~35B's hedged outputs (``Based on the provided context...'') systematically reduce $S_{\text{factuality}}$ below the threshold for clean responses, generating 21--23 false positives per run and causing the system to underperform the naive baseline (71\% vs.\ 85\%). This finding --- the first empirical demonstration of LLM-dependency in NLI-based trust scoring --- implies that the Trust Index threshold $\tau$ and weights must be calibrated per LLM before deployment.

\subsection{RQ3: Security Performance and Operational Overhead}

\textit{How does the inclusion of an Evaluation Agent improve RAG pipeline security, and how does model and retrieval-depth choice affect performance?}

The Evaluation Agent delivers a statistically meaningful security improvement over the naive always-trust baseline when paired with Llama~3.3~70B: 91\% accuracy versus 85\% baseline (+7\%), with zero false positives on clean content. In absolute terms, 6 poisoned contexts out of 15 are correctly identified and blocked before generation, while 9 (the entity-swap and subtle strategy cases) pass the agent undetected.

The 2×2 factorial experiment isolates two key factors:

\textbf{LLM selection is the dominant factor.} Both Llama configurations (MiniLM and Snowflake) produce identical detection results (91\% accuracy, 57.1\% F1), whereas Qwen configurations achieve only 71\% accuracy with significantly lower F1 (31--38\%). The choice of embedding model produces negligible differences within the same LLM. This finding directly implies that system operators choosing models for Trust-Index-based evaluation should prioritise LLM output style over embedding dimensionality or retrieval quality.

\textbf{Retrieval depth does not affect detection.} The K=3 vs.\ K=5 sensitivity analysis (Section~\ref{sec:ch6_topk}) confirms that increasing the number of retrieved documents from 3 to 5 produces no improvement in detection performance for Llama: TP, FP, and FN counts are identical, and the trust score separation changes by only 0.003. The detection ceiling is set by attack strategy difficulty, not by retrieval depth. This null result retroactively validates using K=3 in the 2×2 grid and provides a practical recommendation: K=3 achieves equivalent security at half the NLI inference cost ($\sim$9 vs.\ $\sim$20 calls per batch).

The 13-second per-sample evaluation overhead --- dominated by CPU-based NLI inference and LLM generation latency on the FARMI cluster --- positions the system as suitable for \textbf{batch or high-stakes offline deployment} rather than real-time interactive applications. The evaluation overhead could be reduced substantially through GPU acceleration of the NLI model, which is addressed in Section~\ref{sec:future_work}.

% ─────────────────────────────────────────────────────────────────────────────
\section{Contributions to Trustworthy and Secure AI Research}
\label{sec:disc_contributions}

This thesis makes six distinct contributions spanning theoretical frameworks, empirical findings, and practical system design.

\textbf{Contribution 1: The Trust Index — a unified composite trustworthiness metric.}
Existing evaluation frameworks for RAG systems, such as RAGAS \parencite{es2024ragas} and ARES \parencite{saadfalcon2024ares}, assess whether generated answers are faithful to retrieved context. They do not assess whether the retrieved context itself is trustworthy. The Trust Index ($T = \alpha \cdot S_{\text{factuality}} + \beta \cdot S_{\text{consistency}} + \gamma \cdot (1 - P_{\text{poison}})$) closes this gap by simultaneously quantifying factual grounding, inter-document consistency, and adversarial contamination in a single interpretable score. Its non-linear dampener ensures that strong poisoning signals can override high factuality scores --- a structural safeguard absent from linear aggregation approaches.

\textbf{Contribution 2: A multi-signal Poison Detector for RAG pipelines.}
The five-signal detection architecture (linguistic pattern matching, structural anomaly detection, intra-document NLI, cross-document consistency checking, and semantic outlier detection via embedding centroid deviation) provides defence-in-depth coverage across structurally diverse attack strategies. Relevance-weighted aggregation of per-document poison scores prevents low-relevance suspicious neighbours from dominating the batch-level decision, a design feature that substantially reduced false positives in preliminary evaluations (Round 10 of development, FEVER dataset).

\textbf{Contribution 3: Empirical per-strategy characterisation of detection difficulty.}
The per-strategy evaluation provides the first detailed characterisation of how detection difficulty varies across four adversarial strategies in a RAG-specific context. The hierarchy --- injection (100\% recall) $\gg$ contradiction (33--89\%) $\gg$ entity swap $\approx$ subtle ($\approx$17\%) --- provides concrete guidance for future defence design: strategies that append syntactically distinct content are detectable by surface-level signals; strategies that perform in-place semantic modifications require world-knowledge verification.

\textbf{Contribution 4: Identification of an architectural detection limit.}
Entity swap and subtle manipulation represent a demonstrated \textit{architectural limit} of NLI-and-heuristic approaches: they cannot be reliably detected without access to external factual grounding. This finding advances the scientific understanding of the attack surface in RAG systems and motivates the integration of external knowledge verification (e.g., Wikidata lookup, search-augmented fact checking) in future evaluation agent designs.

\textbf{Contribution 5: First empirical demonstration of LLM-dependency in NLI-based trust scoring.}
While the sensitivity of NLI models to phrasing has been noted in the NLP literature \parencite{honovich2022true}, this thesis provides the first empirical evidence that LLM generation style systematically affects Trust Index values in a RAG evaluation context. The finding that Qwen~3.5~35B's hedged outputs reduce mean clean trust scores by 0.20 points relative to Llama~3.3~70B --- sufficient to push a majority of clean responses below the classification threshold --- has direct implications for deployment: the Trust Index threshold must be treated as an LLM-specific hyperparameter, not a universal constant.

\textbf{Contribution 6: Retrieval-depth invariance finding.}
The K=3 vs.\ K=5 sensitivity analysis establishes a previously unreported property of NLI-based retrieval evaluation: detection performance is insensitive to retrieval depth for the attack strategies and LLM combination tested. This is not obvious a priori (more documents should theoretically improve the chance of exposing poisoned content), and the null result reveals that the bottleneck is the detector's inability to identify certain attack types regardless of batch size. This finding has practical implications for system design: NLI inference cost can be minimised without sacrificing detection quality.

% ─────────────────────────────────────────────────────────────────────────────
\section{Implications for Practical Deployment}
\label{sec:disc_deployment}

The experimental results suggest several concrete guidelines for organisations considering deployment of the Trustworthy RAG framework.

\textbf{Target use cases.} The system's 13-second per-query evaluation overhead makes it unsuitable for real-time conversational applications. Its strongest use case is \textbf{high-stakes batch processing} where response latency is tolerable and the cost of misinformation is high: regulatory document review, medical information retrieval, legal evidence search, and automated fact-checking pipelines. The 100\% recall for instruction injection attacks makes the agent immediately valuable as a defence against prompt injection in document-processing workflows, which represent an increasingly common real-world attack vector \parencite{greshake2023indirect}.

\textbf{LLM selection and per-model calibration.} Operators must calibrate the Trust Index threshold $\tau$ for each LLM used in the system. The fixed $\tau = 0.5$ is empirically tuned for Llama~3.3~70B's concise output style. Verbose or hedged LLMs require a lower threshold (or style-normalisation preprocessing) to avoid systematic false positives on clean content. A practical calibration procedure involves running the evaluation agent on a small held-out set of clean queries for the target LLM and selecting $\tau$ as the 10th percentile of the resulting clean trust score distribution.

\textbf{Strategy-aware deployment.} The per-strategy detection hierarchy implies that the system's value depends heavily on the expected attack distribution. Against adversaries primarily using injection strategies (the most common form of prompt injection), the Evaluation Agent provides near-perfect recall. Against sophisticated adversaries using entity-swap or subtle manipulation techniques, the system's recall falls to approximately 17\%, offering limited additional protection beyond content-agnostic filtering. Operators in adversarial environments should be aware of this coverage gap and supplement the agent with external fact-checking for high-risk queries.

\textbf{Transparency and human oversight.} Each evaluation produces component scores ($S_{\text{factuality}}$, $S_{\text{consistency}}$, $P_{\text{poison}}$), detected warning signals, and a categorical trust level (HIGH / MEDIUM / LOW / VERY LOW), enabling human reviewers to inspect the agent's reasoning rather than relying on a black-box decision. This design aligns with the principle of human-in-the-loop AI systems advocated in the EU AI Act's requirements for high-risk AI applications, particularly those involving safety-critical information provision.

\textbf{Infrastructure recommendations.} GPU acceleration of the BART-large-MNLI model would reduce per-sample evaluation overhead from $\sim$13 seconds to under 2 seconds, enabling near-real-time deployment. Limiting cross-document pairwise NLI to the top-3 retrieved documents (K=3, as validated by the sensitivity analysis) minimises NLI calls without sacrificing detection quality. For systems using hedged LLMs, stripping common hedging prefixes from generated answers before NLI inference would reduce generation-style artefacts without affecting factual content evaluation.

% ─────────────────────────────────────────────────────────────────────────────
\section{Limitations of the Study}
\label{sec:disc_limitations}

This study acknowledges seven principal limitations that bound the generalisability of its findings.

\textbf{Limitation 1: Rule-based adversarial strategies.}
The four poisoning strategies implemented in this thesis are rule-based transformations of existing documents. Real-world adversaries --- particularly those described by \textcite{zou2024poisoned} --- can employ gradient-optimized ``collision'' documents engineered to maximise vector similarity with target queries while containing arbitrary misinformation. Such documents may not exhibit the linguistic artefacts that the current detector relies upon, and their detection would require fundamentally different techniques (e.g., embedding-space anomaly detection or adversarial training of the NLI model).

\textbf{Limitation 2: Detection ceiling for subtle attacks.}
The 16.7\% recall for entity-swap and subtle strategies is a structural limit of the current approach, not a tuning deficiency. Increasing retrieval depth (validated by the K sensitivity analysis), adjusting thresholds, or reweighting the Trust Index components do not improve detection for these strategies. Addressing this limitation requires semantic world-knowledge integration, which is treated as future work.

\textbf{Limitation 3: Evaluation latency.}
The current implementation performs NLI inference on CPU, producing approximately 13 seconds of overhead per sample. This is impractical for interactive deployments. While GPU acceleration would resolve this, the thesis cannot empirically validate the latency improvement due to infrastructure constraints on the FARMI computing cluster.

\textbf{Limitation 4: Fixed classification threshold.}
The threshold $\tau = 0.5$ is constant across query types, domains, and LLMs. It does not adapt to the risk level of the query (a question about medication dosage warrants a lower trust threshold than a general knowledge query), the domain's prior probability of poisoning, or the generation style of the LLM. This inflexibility is a known limitation that the per-LLM calibration recommendation in Section~\ref{sec:disc_deployment} partially addresses but does not fully resolve.

\textbf{Limitation 5: Attack Success Rate not measured.}
The experiments measure whether the Evaluation Agent correctly \textit{flags} poisoned samples, not whether the LLM \textit{incorporates} the misinformation into its answer. True Attack Success Rate measurement would require semantic comparison of generated answers to injected claims. In practice, Llama~3.3~70B frequently ignores appended injection and contradiction content, producing factually correct answers from the legitimate portion of the document. This makes the injected content detectable by the Evaluation Agent even when it fails to influence the generated answer, potentially inflating the perceived utility of detection. Future work should measure ASR directly.

\textbf{Limitation 6: Dataset and domain scope.}
Experiments use two standardised benchmarks (TruthfulQA and FEVER) under controlled poisoning conditions. Performance on specialised corpora --- medical literature, legal documents, scientific papers --- may differ substantially, as these domains involve technical vocabulary, complex claim structures, and different LLM behaviour patterns. The generalisability of the Trust Index to such domains has not been empirically validated.

\textbf{Limitation 7: Ground-truth convention for poisoning labels.}
The current evaluation protocol labels a sample as poisoned based on document-level index alignment (whether the knowledge-base document corresponding to a query was modified), rather than directly checking whether a poisoned document was retrieved in that specific query's top-$K$ results. This convention is suitable for the controlled question-document paired benchmark setup used in this thesis, but it may overestimate or underestimate detection metrics in settings where retrieval frequently returns off-index neighbours. Future evaluations should implement retrieval-grounded poisoning labels to improve ecological validity.

% ─────────────────────────────────────────────────────────────────────────────
\section{Alignment with Ethical AI and Security Frameworks}
\label{sec:disc_ethics}

The development and evaluation of adversarial AI systems carries inherent ethical responsibilities. This section reflects on the ethical dimensions of the thesis in the context of established AI governance frameworks.

\textbf{Transparency and explainability.} The Evaluation Agent is explicitly designed to produce interpretable outputs: component trust scores, detected warning signals, and a categorical trust level accompany every evaluation decision. This design aligns with the transparency and explainability requirements of the EU AI Act (Article 13) for high-risk AI systems, which mandate that automated decisions be explainable to affected parties. Unlike black-box LLM-as-a-Judge approaches, the agent's decisions can be audited and challenged by human reviewers.

\textbf{Defensive intent and responsible disclosure.} The adversarial dataset generator and poisoning strategies implemented in this thesis are designed exclusively for controlled research evaluation. The code does not include optimized attack capabilities (such as gradient-based collision document generation) that would provide operational uplift to malicious actors. The strategies implemented are simplified, rule-based transformations intended to characterise detection difficulty, not to enable novel attacks. This approach aligns with responsible disclosure principles in security research: the system's limitations (particularly the near-zero recall for entity-swap attacks) are transparently documented, enabling downstream researchers and practitioners to make informed deployment decisions.

\textbf{Human-in-the-loop design.} The Evaluation Agent is designed to augment rather than replace human oversight. It provides trust scores and warnings that inform human decision-making rather than autonomously filtering or censoring retrieved content. The categorical trust levels (HIGH / MEDIUM / LOW / VERY LOW) are calibrated to support graduated human review: HIGH-trust responses may be used directly, while VERY LOW-trust responses should be independently verified before use. This design philosophy aligns with the principle of meaningful human control advocated by the OECD Principles on Artificial Intelligence and the IEEE Ethically Aligned Design framework.

\textbf{Dual-use considerations.} The detection techniques developed in this thesis --- particularly the linguistic pattern library and intra-document NLI --- could theoretically be inverted to inform adversaries about the specific features that the detector monitors, enabling more sophisticated evasion. This dual-use risk is inherent to all security research and is mitigated by the open publication of both the detection system and its known limitations, consistent with the principle that defensive knowledge should be equally accessible to defenders.

% ─────────────────────────────────────────────────────────────────────────────
\section{Summary}
This chapter interpreted the experimental results in the context of the three Research Questions, articulated six distinct contributions, and identified limitations and deployment implications. The core finding is that the Evaluation Agent provides reliable, interpretable defence against overt poisoning strategies (instruction injection and contradiction) at the cost of an evaluation overhead that currently limits it to batch deployments. Entity-swap and subtle manipulation represent a fundamental architectural limit requiring world-knowledge integration to overcome. The LLM-dependency of the Trust Index and the retrieval-depth-invariance finding are novel empirical contributions that advance the design science of trustworthy RAG evaluation. The following chapter synthesises these findings into a final conclusion and outlines a research roadmap toward secure, auditable RAG systems.

\chapter{Conclusion and Future Work}
\label{ch8_conclusion}

This thesis investigated the design and evaluation of an autonomous Evaluation Agent for detecting misinformation and knowledge poisoning in Retrieval-Augmented Generation (RAG) systems. Motivated by the growing deployment of RAG in high-stakes information environments and the demonstrable vulnerability of retrieval-augmented architectures to corpus poisoning attacks, the work addressed a gap left by existing evaluation frameworks: the absence of a runtime defence that simultaneously assesses factual grounding, inter-document consistency, and adversarial contamination.

% ─────────────────────────────────────────────────────────────────────────────
\section{Summary of Research Objectives and Findings}
\label{sec:conc_summary}

The thesis was structured around three Research Questions, each of which is now answered by the empirical evidence gathered across the experimental programme.

\textbf{RQ1 — Detection Effectiveness.}
\textit{How effectively can an autonomous Evaluation Agent detect misinformation and knowledge poisoning in a RAG pipeline?}

The Evaluation Agent achieves high effectiveness against overt poisoning strategies but encounters a structural detection limit against subtle in-place modifications. Instruction injection is detected with 100\% recall; contradiction poisoning achieves 33--89\% recall depending on experimental conditions; entity-swap and subtle manipulation are detected at approximately 17\% recall --- a limit imposed by the absence of external world-knowledge reasoning in the current architecture. In the primary mixed-strategy 100-sample experiment, the system achieves 40\% overall recall with 100\% precision (zero false positives on clean content), confirming that the agent's detections are highly reliable when they occur.

\textbf{RQ2 — Trust Index Reliability.}
\textit{How reliably does the Trust Index quantify trustworthiness across diverse poisoning scenarios?}

The Trust Index formula $T = 0.4 \cdot S_{\text{factuality}} + 0.35 \cdot S_{\text{consistency}} + 0.25 \cdot (1 - P_{\text{poison}})$, augmented with a non-linear dampener for high-contamination conditions, provides a calibrated and discriminative trustworthiness score. Mean trust score separation between clean (0.830) and poisoned (0.590) contexts confirms meaningful discriminative power. The ablation study validates the default weight configuration as the best overall trade-off. However, the 2×2 factorial experiment reveals that Trust Index values are LLM-dependent: verbose generation styles systematically reduce $S_{\text{factuality}}$ for clean responses, requiring per-LLM threshold calibration before production deployment.

\textbf{RQ3 — Security Improvement and Overhead.}
\textit{How does the Evaluation Agent improve RAG security compared to a baseline, and how do model and retrieval-depth choices affect performance?}

The Evaluation Agent improves overall accuracy by +7\% over the naive always-trust baseline (91\% vs.\ 85\%) with Llama~3.3~70B and zero false positives on clean content, providing a +5\% improvement on the FEVER benchmark. The 2×2 factorial experiment establishes that LLM selection is the dominant performance factor: both Llama embedding configurations achieve identical results, while Qwen~3.5~35B underperforms the baseline by 16\% due to generation-style false positives. The K=3 vs.\ K=5 sensitivity analysis confirms that detection performance is retrieval-depth-invariant, and that K=3 is the cost-optimal retrieval depth for this system.

Table~\ref{tab:conc_rq_summary} summarises the correspondence between each Research Question and its primary experimental evidence.

\begin{table}[h]
\centering
\caption{Research Questions and corresponding experimental findings.}
\label{tab:conc_rq_summary}
\begin{tabular}{lp{5cm}p{6cm}}
\hline
\textbf{RQ} & \textbf{Finding} & \textbf{Evidence} \\
\hline
RQ1 & High recall for injection/contradiction; near-zero for entity-swap/subtle & Per-strategy experiments (Section~\ref{sec:ch6_detection}) \\
RQ2 & Trust Index discriminates clean vs.\ poisoned; LLM-dependent & Main experiment + 2×2 grid (Sections~\ref{sec:ch6_baseline}, \ref{sec:ch6_grid}) \\
RQ3 & +7\% vs baseline; LLM $\gg$ embedding effect; K-invariant & 2×2 grid + K sensitivity (Sections~\ref{sec:ch6_grid}, \ref{sec:ch6_topk}) \\
\hline
\end{tabular}
\end{table}

% ─────────────────────────────────────────────────────────────────────────────
\section{Theoretical and Practical Contributions}
\label{sec:conc_contributions}

\subsection{Theoretical Contributions}

\textbf{The Trust Index as a unified security-reliability metric.}
This thesis formalises the first composite metric designed to simultaneously evaluate the factual grounding, internal consistency, and adversarial safety of a RAG response in a single interpretable score. Existing frameworks such as RAGAS \parencite{es2024ragas} and ARES \parencite{saadfalcon2024ares} address faithfulness but not corpus integrity. The Trust Index bridges this gap and establishes a mathematical foundation for what the thesis terms the ``Trust-Security Gap'' in RAG evaluation.

\textbf{Empirical characterisation of per-strategy detection difficulty.}
The per-strategy experimental analysis provides empirical evidence of the varying detectability of four structurally distinct adversarial strategies in RAG corpora. The resulting detection hierarchy (injection $\gg$ contradiction $\gg$ entity-swap $\approx$ subtle) constitutes a reusable reference for future research on RAG defence design, establishing which attack families are amenable to surface-signal detection and which require semantic world-knowledge integration.

\textbf{LLM-dependency of NLI-based trust scoring.}
The 2×2 factorial experiment produces the first systematic empirical evidence that NLI-based trust assessment is sensitive to LLM generation style. This finding generalises beyond the specific models tested: any LLM whose output distribution differs significantly from the Trust Index's implicit calibration target will require per-model threshold adjustment, a consideration absent from prior Trust Index proposals.

\textbf{Retrieval-depth invariance in NLI-based evaluation.}
The K sensitivity analysis demonstrates a non-obvious architectural property: detection performance in NLI-based retrieval evaluation is insensitive to retrieval depth for the tested attack strategies. The null result (K=5 $\equiv$ K=3) advances the theoretical understanding of where evaluation agent performance is and is not bounded by retrieval quality.

\subsection{Practical Contributions}

\textbf{A working, reproducible implementation.}
The complete Trustworthy RAG framework is implemented in Python 3.10 using LangChain, FAISS, HuggingFace Transformers (BART-large-MNLI), and an OpenAI-compatible API for LLM access. The implementation is modular and documented, with independently replaceable components for the retriever, NLI verifier, poison detector, and trust calculator.

\textbf{Adversarial dataset generator.}
The \textit{PoisonedDatasetGenerator} implements four rule-based poisoning strategies with configurable poison ratios and supports both TruthfulQA and FEVER benchmarks. It produces reproducible experimental datasets for future comparative studies and is extensible to additional strategies.

\textbf{Reproducible experiment framework.}
The experiment runner supports automated multi-model grid evaluation ($2 \times 2$ factorial design), per-strategy isolated experiments, ablation studies, and retrieval depth sensitivity analysis. Results are automatically serialised to JSON for offline analysis and reproducibility verification.

% ─────────────────────────────────────────────────────────────────────────────
\section{Recommendations for Future Research}
\label{sec:future_work}

The findings of this thesis identify seven high-priority directions for future research.

\textbf{1. Gradient-optimized adversarial evaluation.}
The four strategies evaluated in this thesis are rule-based approximations of real adversarial attacks. Future work should evaluate the Evaluation Agent against gradient-optimized collision documents of the type described by \textcite{zou2024poisoned}, which are engineered to maximise retrieval similarity while containing arbitrary misinformation. These attacks are likely to evade linguistic pattern detection and may require NLI-adversarial training or embedding-space anomaly detection to identify.

\textbf{2. Per-LLM Trust Index calibration.}
The LLM-dependency finding motivates the development of a systematic calibration procedure for the Trust Index threshold $\tau$ and component weights $\alpha$, $\beta$, $\gamma$. A practical approach would use a small LLM-specific validation set of clean queries to fit $\tau$ to the 10th percentile of the clean trust score distribution, or to train a lightweight per-LLM calibration layer that normalises NLI entailment scores for generation-style artefacts.

\textbf{3. Semantic world-knowledge integration.}
Entity-swap and subtle manipulation represent the architectural limit of the current surface-signal approach. Integrating external knowledge verification --- for example, querying structured knowledge bases such as Wikidata to verify entity-attribute relationships, or employing a search-augmented fact checker for open-domain claims --- could substantially improve recall for these attack families. The key design challenge is maintaining acceptable latency while adding an external verification call to each retrieved document.

\textbf{4. GPU-accelerated evaluation pipeline.}
Deploying the NLI model on GPU hardware would reduce per-sample evaluation latency from approximately 13 seconds to below 2 seconds, enabling near-real-time RAG evaluation. The current CPU-only implementation is a practical constraint of the FARMI cluster's resource allocation for student projects, not an inherent architectural limitation.

\textbf{5. True Attack Success Rate measurement.}
The current experimental framework measures the Evaluation Agent's \textit{detection} of poisoned samples, not the LLM's \textit{adoption} of injected misinformation. Future work should implement ASR measurement via embedding-similarity comparison between generated answers and injected claims, providing a complete end-to-end evaluation of the system's security value.

\textbf{6. Retrieval-grounded poisoning labels.}
The current benchmark protocol uses document-level index-aligned poisoning labels to determine whether each sample is poisoned. This is appropriate for the controlled paired benchmark setup used in this thesis, but future work should label poisoning based on whether at least one poisoned document is actually retrieved in the query's top-$K$ set. This would strengthen external validity in open-corpus settings where neighbour retrieval patterns are less deterministic.

\textbf{7. Multimodal and multi-domain extension.}
This thesis evaluates the framework on text-only English benchmarks (TruthfulQA and FEVER). Future research should extend the evaluation to domain-specific corpora (medical, legal, scientific), multilingual settings, and multimodal RAG systems incorporating images and structured data, where both the poisoning attack surface and the detection approach may differ substantially.

% ─────────────────────────────────────────────────────────────────────────────
\section{Roadmap Toward Secure and Auditable RAG Systems}
\label{sec:conc_roadmap}

The Trustworthy RAG framework developed in this thesis represents a first step toward a broader vision: RAG systems that are not merely accurate but \textit{auditable} --- systems that provide verifiable evidence for their trustworthiness claims alongside each generated response.

\textbf{Near-term (1--2 years):} The most impactful immediate improvements are GPU-based NLI inference for real-time operation, per-LLM calibration procedures for Trust Index thresholds, and extension to a broader set of LLMs and embedding models. Standardising the Trust Index as a reporting metric in RAG evaluation benchmarks (alongside existing faithfulness and relevance metrics) would help establish it as a community baseline.

\textbf{Medium-term (2--4 years):} Integrating external knowledge verification APIs (structured databases, search engines, specialised fact-checking services) into the evaluation pipeline would directly address the detection ceiling for in-place modification attacks. Streaming evaluation architectures --- where the Evaluation Agent operates on individual retrieved passages as they are returned by the retriever rather than on the full batch --- would reduce latency and support real-time interactive deployments.

\textbf{Long-term (4+ years):} The ultimate goal is \textit{multi-agent evaluation}: architectures in which specialised agents handle factuality verification, security assessment, consistency checking, and provenance tracking as independent, composable modules. Combined with tamper-evident audit logs and cryptographic provenance chains for retrieved documents, such systems would provide end-to-end accountability for AI-generated information --- a critical requirement for deploying generative AI in regulated environments such as healthcare, finance, and public administration.

The vulnerability of RAG pipelines to knowledge poisoning is not a temporary implementation flaw but a structural consequence of the architecture's reliance on unverified external content. As RAG systems become increasingly central to knowledge-intensive applications, the need for systematic, runtime trustworthiness evaluation will only grow. This thesis establishes that such evaluation is technically feasible, provides a validated framework and implementation, and identifies the specific challenges that must be overcome to make it universally deployable.

% ─────────────────────────────────────────────────────────────────────────────
\section{Closing Remark}

Knowledge Poisoning turns the core strength of RAG systems --- their ability to ground responses in external evidence --- into a liability. By demonstrating that an autonomous Evaluation Agent can reliably detect overt poisoning strategies, characterise the architectural limits of surface-signal detection, and produce interpretable trust scores with validated discriminative power, this thesis advances the science of trustworthy generative AI. The path from the current system to one that is universally secure, real-time, and domain-agnostic is long, but the foundation --- a principled, empirically validated Trust Index and a multi-signal defence architecture --- is now established.


% ---------- References ----------
\printbibliography

% ---------- Appendices ----------
\appendix
\chapter{System Configuration Parameters}
\label{app:parameters}

Table~\ref{tab:app_parameters} lists the complete set of configurable parameters used in the Trustworthy RAG framework, with default values and valid ranges.

\begin{table}[h]
\centering
\caption{System configuration parameters and default values.}
\label{tab:app_parameters}
\begin{tabularx}{\textwidth}{@{}llcX@{}}
\hline
\textbf{Parameter} & \textbf{Default} & \textbf{Range} & \textbf{Description} \\
\hline
\texttt{TOP\_K\_RETRIEVAL} & 5 & 1--10 & Number of documents retrieved per query \\
\texttt{TRUST\_THRESHOLD} ($\tau$) & 0.5 & 0.0--1.0 & Binary trustworthiness decision boundary \\
\texttt{POISON\_RATIO} & 0.30 & 0.0--1.0 & Fraction of poisoned documents in the corpus \\
$\alpha$ (factuality weight) & 0.40 & 0.0--1.0 & Trust Index weight for NLI factuality score \\
$\beta$ (consistency weight) & 0.35 & 0.0--1.0 & Trust Index weight for inter-document consistency \\
$\gamma$ (poison weight) & 0.25 & 0.0--1.0 & Trust Index weight for poison safety component \\
Dampener threshold & 0.70 & 0.5--1.0 & $P_{\text{poison}}$ value above which the non-linear dampener activates \\
Dampener max reduction & 0.40 & 0.0--1.0 & Maximum multiplicative trust score reduction at $P_{\text{poison}}=1.0$ \\
FAISS confidence floor & 0.30 & 0.0--1.0 & Minimum retrieval similarity score before a trust penalty is applied \\
NLI entailment threshold & 0.40 & 0.0--1.0 & Minimum entailment score for a document to contribute to factuality \\
Cross-doc contradiction threshold & 0.70 & 0.5--1.0 & Threshold for high-severity cross-document contradiction signal \\
Intra-doc contradiction threshold & 0.70 & 0.5--1.0 & Threshold for intra-document first/second half NLI contradiction \\
\texttt{NLI\_MODEL} & \multicolumn{2}{c}{\texttt{facebook/bart-large-mnli}} & NLI model used for factuality and poison detection \\
\texttt{EMBEDDING\_MODEL} & \multicolumn{2}{c}{\texttt{all-MiniLM-L6-v2}} & Sentence embedding model for vector retrieval \\
\texttt{LLM\_MODEL} & \multicolumn{2}{c}{\texttt{llama3.3:70b}} & Primary LLM used for response generation \\
\hline
\end{tabularx}
\end{table}

\chapter{Extended Experimental Results}
\label{app:results}

This appendix provides supplementary results tables referenced in Chapter~\ref{ch6_experiments}.

\section{Per-Strategy Results — All Metrics}

Table~\ref{tab:app_perstrategy} extends the per-strategy summary with full TP/FP/FN counts and trust score separation values (100 samples per strategy, TruthfulQA, $K=5$).

\begin{table}[h]
\centering
\caption{Per-strategy detection results (Llama~3.3~70B + all-MiniLM-L6-v2, 100 samples each, TruthfulQA, $K=5$). Results from isolated single-strategy experiments; mixed-strategy results reported in Chapter~\ref{ch6_experiments}.}
\label{tab:app_perstrategy}
\begin{tabular}{lcccccccc}
\hline
\textbf{Strategy} & \textbf{Acc.} & \textbf{Prec.} & \textbf{Recall} & \textbf{F1} & \textbf{Sep.} & \textbf{TP} & \textbf{FP} & \textbf{FN} \\
\hline
Contradiction        & 92\% & 88.9\% & 53.3\% & 66.7\% & 0.311 & 8  & 1  & 7  \\
Instruction Injection & 99\% & 93.8\% & 100\%  & 96.8\% & 0.498 & 15 & 1  & 0  \\
Entity Swap          & 85\% & ---    & 0\%    & 0\%    & 0.053 & 0  & 0  & 15 \\
Subtle               & 88\% & 100\%  & 20.0\% & 33.3\% & 0.149 & 3  & 0  & 12 \\
Mixed                & 91\% & 100\%  & 40.0\% & 57.1\% & 0.225 & 6  & 0  & 9  \\
\hline
\end{tabular}
\end{table}

\section{Ablation Study Results}

Table~\ref{tab:app_ablation} reports the full ablation study results across five Trust Index weight configurations.

\begin{table}[h]
\centering
\caption{Ablation study: Trust Index weight configurations and detection performance (TruthfulQA, 60 samples).}
\label{tab:app_ablation}
\begin{tabular}{lccccccc}
\hline
\textbf{Configuration} & $\alpha$ & $\beta$ & $\gamma$ & \textbf{Acc.} & \textbf{Prec.} & \textbf{Recall} & \textbf{F1} \\
\hline
Default             & 0.40 & 0.35 & 0.25 & 86.7\% & 100\% & 11.1\% & 20.0\% \\
Equal weights       & 0.33 & 0.34 & 0.33 & 86.7\% & 100\% & 11.1\% & 20.0\% \\
Factuality-heavy    & 0.60 & 0.20 & 0.20 & 86.7\% & 100\% & 11.1\% & 20.0\% \\
Consistency-heavy   & 0.20 & 0.60 & 0.20 & 88.3\% & 100\% & 22.2\% & 36.4\% \\
Poison-heavy        & 0.20 & 0.20 & 0.60 & 88.3\% & 100\% & 22.2\% & 36.4\% \\
\hline
\end{tabular}
\end{table}

\section{FEVER Benchmark Pilot Results}

Table~\ref{tab:app_fever} reports early pilot results on the FEVER benchmark (Llama~3.3~70B + all-MiniLM-L6-v2). These 20-sample pilot results indicated high accuracy (90\%) initially, but a subsequent full 100-sample experiment (reported in Chapter~\ref{ch6_experiments}) revealed broader generalisability challenges.

\begin{table}[h]
\centering
\caption{FEVER benchmark early pilot results after KB enrichment fix (20 samples, mixed strategy, $K=3$).}
\label{tab:app_fever}
\begin{tabular}{lcccccc}
\hline
\textbf{Metric} & \textbf{Value} & \quad & \textbf{Metric} & \textbf{Value} \\
\hline
Accuracy        & 90\%   & & Clean Trust $\bar{T}$ & 0.654 \\
Precision       & ---    & & Poison Trust $\bar{T}$ & --- \\
Recall          & 66.7\% & & Separation    & --- \\
F1              & ---    & & vs.\ Baseline & +5\% \\
\hline
\end{tabular}
\end{table}

\chapter{LLM Prompt Templates}
\label{app:prompts}

This appendix lists the prompt templates used in the RAG pipeline for response generation and the system prompt governing the LLM's behaviour.

\section{RAG Generation Prompt}

The following template is used by the generator to produce responses from retrieved context. The placeholders \texttt{\{context\}} and \texttt{\{question\}} are filled at runtime by the pipeline orchestrator.

\begin{verbatim}
You are a helpful assistant. Answer the question based ONLY on the
provided context. If the context does not contain enough information
to answer the question, say "I don't have enough information to
answer this question."

Context:
{context}

Question: {question}

Answer:
\end{verbatim}

\section{Design Rationale}

The prompt is intentionally minimal and instruction-focused. The phrase ``based ONLY on the provided context'' is included to maximise the LLM's reliance on retrieved documents (promoting faithfulness) and to reduce parametric knowledge injection (which would undermine the retrieval-grounded evaluation). The explicit fallback instruction (``I don't have enough information...'') prevents the model from hallucinating answers when retrieved context is insufficient, which would otherwise produce low NLI entailment scores and artificially depressed Trust Index values.

No chain-of-thought or few-shot examples are included, consistent with the evaluation design principle that the generation step should reflect a standard, unaugmented RAG deployment rather than a reasoning-enhanced configuration.

\chapter{Ethical Considerations and Responsible Research Statement}
\label{app:ethics}

\section{Research Ethics Declaration}

This thesis conducts research on adversarial attacks and defences in AI systems. All experiments were performed in a controlled, isolated research environment using publicly available benchmark datasets (TruthfulQA and FEVER). No personally identifiable information was collected, processed, or stored at any stage of the research.

The adversarial dataset generator and poisoning strategies implemented in this thesis are designed exclusively for the purpose of evaluating the Evaluation Agent's detection capabilities. They are not intended to serve as operational attack tools. The strategies are rule-based transformations (not gradient-optimized collision attacks) and do not provide meaningful uplift to malicious actors beyond what is already published in the open literature \parencite{zou2024poisoned, greshake2023indirect}.

\section{Dual-Use Considerations}

The detection techniques developed in this thesis --- particularly the linguistic pattern library and the intra-document NLI approach --- could theoretically inform adversaries about the specific features monitored by the detector. This dual-use risk is inherent to all published security research and is mitigated by the transparent documentation of both the detection system's capabilities and its known limitations (particularly the near-zero recall for entity-swap and subtle manipulation strategies). The publication of defensive limitations serves the broader research community by directing future work toward the correct problem: detection of in-place semantic modifications using external world-knowledge, rather than surface-level textual signals.

\section{Data Handling}

All datasets used in this research are publicly licensed:
\begin{itemize}
    \item \textbf{TruthfulQA} \parencite{lin2022truthfulqa}: Apache License 2.0
    \item \textbf{FEVER} \parencite{thorne2018fever}: Creative Commons Attribution-ShareAlike 3.0 Unported
\end{itemize}

The synthetic poisoned datasets generated during experiments are derived from these sources and carry the same license restrictions. No proprietary, confidential, or sensitive data sources were used.

\section{AI-Assisted Research Tools}

Portions of the implementation code and thesis writing process benefited from AI-assisted development tools (Claude Code, Anthropic). All AI-generated content was reviewed, validated, and edited by the thesis author. The experimental results, analysis, and conclusions are the independent work of the author, verified against the actual codebase and experiment outputs.

\chapter{Implementation Details and Repository Structure}
\label{app:implementation}

\section{Repository Structure}

The project repository is organised as follows:

\begin{verbatim}
RAG Agent/
├── src/
│   ├── retriever/
│   │   ├── embeddings.py        # Embedding model wrappers (local + API)
│   │   └── retriever.py         # FAISS-based document retriever
│   ├── evaluation_agent/
│   │   ├── nli_verifier.py      # BART-MNLI NLI inference
│   │   ├── poison_detector.py   # Five-signal poison detection
│   │   ├── trust_index.py       # Trust Index computation + dampener
│   │   └── evaluation_agent.py  # Orchestration class
│   ├── generator/
│   │   └── llm_generator.py     # OpenAI-compatible LLM interface
│   └── pipeline/
│       ├── rag_pipeline.py      # End-to-end RAG pipeline
│       └── experiment_runner.py # Experiment execution + metrics
├── src/experiments/
│   └── poisoned_dataset.py      # Adversarial dataset generator
├── configs/
│   └── config.yaml              # Master configuration file
├── data/
│   ├── raw/                     # Downloaded benchmark datasets
│   ├── processed/               # Preprocessed knowledge bases
│   └── experiments/             # Serialised experiment results (JSON)
├── Master_Thesis/
│   ├── main.tex                 # LaTeX document root
│   └── chapters/                # Chapter and appendix .tex files
├── run_experiment.py            # CLI experiment runner
└── requirements.txt             # Python package dependencies
\end{verbatim}

\section{Key Dependencies}

\begin{table}[h]
\centering
\caption{Core Python dependencies and versions.}
\label{tab:app_deps}
\begin{tabular}{lll}
\hline
\textbf{Package} & \textbf{Version} & \textbf{Purpose} \\
\hline
Python         & 3.10+    & Runtime language \\
LangChain      & 0.1+     & Pipeline orchestration \\
faiss-cpu      & 1.7+     & Vector similarity search \\
transformers   & 4.35+    & BART-MNLI NLI model \\
sentence-transformers & 2.2+ & MiniLM-L6-v2 embeddings \\
torch          & 2.0+     & NLI model inference (CPU) \\
openai         & 1.0+     & FARMI API client \\
datasets       & 2.14+    & HuggingFace dataset loading \\
loguru         & 0.7+     & Structured logging \\
\hline
\end{tabular}
\end{table}

\section{Reproduction Instructions}

To reproduce the primary experiment results (Llama~3.3~70B + all-MiniLM-L6-v2, TruthfulQA, 100 samples, mixed strategy, K=3):

\begin{enumerate}
    \item Install dependencies: \texttt{pip install -r requirements.txt}
    \item Configure FARMI API credentials in \texttt{.env}: set \texttt{FARMI\_API\_KEY} and \texttt{FARMI\_API\_BASE}
    \item Run: \texttt{python run\_experiment.py}; at the interactive prompt, select Llama~3.3~70B, all-MiniLM-L6-v2, and K=3
    \item Results are saved to \texttt{data/experiments/} as a timestamped JSON file
\end{enumerate}

To reproduce the full 2×2 grid (K=3):

\begin{verbatim}
python run_experiment.py --grid --samples 50
\end{verbatim}

To reproduce the K=5 sensitivity analysis:

\begin{verbatim}
python run_experiment.py --grid --samples 50
# Select K=5 at the interactive prompt, or set TOP_K_RETRIEVAL=5 in .env
\end{verbatim}

All experiment results used in this thesis are archived in \texttt{data/experiments/} with ISO 8601 timestamps in the filename. The primary results file is:\\
\texttt{experiment\_truthfulqa\_llama3.3-70b\_all-MiniLM-L6-v2\_detection\_20260307\_112613.json}


\end{document}
